\documentclass[11pt,a4paper]{article}
\usepackage[margin=1.8cm]{geometry}
\usepackage{amsmath,amssymb,tabularx,array,needspace,fontspec,cancel}
\setmainfont{FreeSerif}
\newfontfamily\urdufont[Script=Arabic]{FreeSerif}
\TeXXeTstate=1
\newcommand{\urdu}[1]{{\urdufont\beginR #1\endR}}
\newcommand{\lr}[1]{\beginL #1\endL}
\setlength{\parindent}{0pt}
\setlength{\parskip}{0pt}
\renewcommand{\arraystretch}{1.15}
\newcommand{\tb}{\cancel{||||}\,}
% \umcq{n}{urdu question}{A}{B}{C}{D}; options shown D C B A left-to-right like the original
\newcommand{\umcq}[6]{%
  \Needspace{4\baselineskip}%
  \par\medskip{\urdufont\raggedleft\beginR\textbf{\lr{#1.}}\ #2\endR\par}\nopagebreak
  \noindent\begin{tabularx}{\linewidth}{@{}*{4}{>{\raggedleft\arraybackslash}X}@{}}
    \lr{(D)}~#6 & \lr{(C)}~#5 & \lr{(B)}~#4 & \lr{(A)}~#3
  \end{tabularx}\par}

\begin{document}

\begin{center}
  {\Huge\bfseries The Hope Science Academy}\\[3pt]
  {\small Excellence is our standard not goal}\\[2pt]
  {\small 247-E-1, (Opposite LDA School) Johar Town, Lahore. PH: 0300-8194789, 0322-4157001}
\end{center}

\medskip
\noindent
\begin{tabularx}{\linewidth}{|X|X|X|X|}
\hline
\textbf{Student Name:} & \textbf{Roll Num:} & \textbf{Class:} 9th & \textbf{Paper Code:} \\ \hline
\textbf{Subject Name:} Mathematics. & \textbf{Time Allowed:} & \textbf{Total Marks:} 543 & \textbf{Exam Date:} 09-Oct-2026 \\ \hline
\multicolumn{4}{|l|}{\textbf{Exam Syllabus:} FULL BOOK} \\ \hline
\end{tabularx}

\bigskip
\noindent\textbf{1$\times$543=543}\hfill{\urdufont\beginR\textbf{سوال نمبر \lr{1}۔ چار ممکنہ جوابات میں سے درست پر دائرہ لگائیں}\endR}

\umcq{1}{سمیریوں نے کون سا عددی نظام استعمال کیا:}{\urdu{اساس}}{\urdu{اساس}}{\urdu{اساس}}{\urdu{اساس}}
\umcq{2}{صفر کا تصور سب سے پہلے تیار کیا گیا تھا:}{\urdu{مصریوں}}{\urdu{سمیریوں}}{\urdu{انڈین}}{\urdu{یونانی}}
\umcq{3}{جدید عددی نظام جیسے ہیکسا ڈیسیملکس اساس پر ہیں:}{\urdu{اساس}}{\urdu{اساس}}{\urdu{اساس}}{\urdu{اساس}}
\umcq{4}{درج ذیل میں سے ناطق عدد ہے:}{0.333...}{3.14159}{1.5}{2332}
\umcq{5}{متوالی کسور اعشاریہ ہمیشہ ہوتی ہے:}{\urdu{ناطق عدد}}{\urdu{غیر ناطق عدد}}{\urdu{عدد}}{\urdu{مکمل عدد}}
\umcq{6}{کون سا عدد ناطق ہے:}{$\sqrt{3}$}{$\sqrt{16}$}{$\sqrt{5}$}{$\sqrt{7}$}
\umcq{7}{\lr{$\dfrac{1}{7}$} کا اعشاری پھیلاؤ ہے::}{\urdu{مختتم کسور اعشاریہ}}{\urdu{غیر مختتم اور غیر متوالی کسر اعشاریہ}}{\urdu{غیر مختتم اور متوالی کسر اعشاریہ}}{\urdu{ان میں سے کوئی نہیں}}
\umcq{8}{\lr{a+b = b + a} سے کون سی خاصیت بیان ہوتی ہے:}{\urdu{خاصیت تلازم}}{\urdu{جمعی ذاتی عنصر}}{\urdu{تقسیمی}}{\urdu{خاصیت مبادلہ}}
\umcq{9}{حقیقی اعداد کی جمع خاصیت ہے:}{1}{0}{-1}{None of these}
\umcq{10}{\lr{a(b+c)=ab+ac} کونسی خاصیت کو بیان کرتا ہے:}{\urdu{خاصیت تلازم}}{\urdu{تقسیمی}}{\urdu{ضربی معکوس}}{\urdu{خاصیت بندش}}
\umcq{11}{\lr{5} کو ضربی معکوس ہے:}{-5}{$\dfrac{1}{5}$}{0}{1}
\umcq{12}{اگر \lr{$a>b$} اور \lr{$b>c$}، ہو تو \lr{$a>c$} ...... خاصیت کہلاتی ہے:}{\urdu{خاصیت عکسی}}{\urdu{خاصیت تشاکل}}{\urdu{خاصیت متعدیت}}{\urdu{خاصیت ثلاثی}}
\umcq{13}{مندرجہ ذیل میں سے کون سا عدد نہیں ہے:}{-3}{0}{4.5}{7}
\umcq{14}{درج ذیل میں سے کون سا ناطق اعداد کی مثال ہے:}{$\pi$}{$\dfrac{3}{4}$}{$\sqrt{2}$}{e}
\umcq{15}{متوالی کسر اعشاریہ کی مثال کیا ہے:}{3.14159265 ...}{0.3333 ...}{2.718281828 ...}{0.49835555}
\umcq{16}{\lr{-3} کون سا نمبر سسٹم سے تعلق رکھتا ہے:}{\urdu{قدرتی اعداد}}{\urdu{مکمل اعداد}}{\urdu{اعداد}}{\urdu{غیر ناطق اعداد}}
\umcq{17}{درج ذیل میں سے کون سا حقیقی عدد نہیں ہے:}{5}{-2}{$\sqrt{7}$}{0}
\umcq{18}{\lr{0} کس نمبر سیٹ سے تعلق رکھتا ہے:}{\urdu{مکمل اعداد}}{\urdu{قدرتی اعداد}}{\urdu{غیر ناطق اعداد}}{\urdu{ان میں سے کوئی نہیں}}
\umcq{19}{مندرجہ ذیل میں سے کون سا لامحدود غیر متوالی کسر اعشاریہ کی مثال ہے:}{0.25}{0.333...}{5.75}{$\pi$}
\umcq{20}{\lr{$0.\overline{6}$} کس کے برابر ہوتا ہے:}{$\dfrac{2}{3}$}{$\dfrac{1}{6}$}{$\dfrac{3}{5}$}{$\dfrac{6}{99}$}
\umcq{21}{کون سا دو رقمی مقادیر اصم ہے:}{$\sqrt{3}$}{$\sqrt{3}+\sqrt{5}$}{$2\sqrt{3}$}{$\sqrt{9}$}
\umcq{22}{\lr{$\sqrt{7}-\sqrt{2}$} کا کونجوگیٹ ہے::}{$\sqrt{7}+\sqrt{2}$}{$\sqrt{7}-\sqrt{2}$}{$-\sqrt{7}+\sqrt{2}$}{$-\sqrt{7}-\sqrt{2}$}
\umcq{23}{مخرج کا ناطق ہے: \lr{$\dfrac{1}{\sqrt{5}}$} :}{$\sqrt{5}$}{$\dfrac{\sqrt{5}}{5}$}{$\dfrac{1}{5}$}{$5\sqrt{5}$}

\umcq{24}{\lr{:}}{$\sqrt{8}$}{$\sqrt{12}$}{$\sqrt{16}$}{$\sqrt{18}$}
\umcq{25}{\lr{2x+5=15} کا حل کیا ہوگا:}{x=3}{x=10}{x=7.5}{x=5}
\umcq{26}{اگر \lr{$\sqrt{x}=5$} ہو تو \lr{x=?}:}{10}{25}{5}{15}
\umcq{27}{\lr{$\dfrac{\sqrt{12}\times\sqrt{3}}{\sqrt{4}}$} کی قیمت کیا ہوگی::}{3}{6}{9}{12}
\umcq{28}{اگر \lr{a=3,b=4} ہو تو \lr{$\sqrt{a^{2}+b^{2}}$}}{5}{7}{12}{25}
\umcq{29}{\lr{$\dfrac{\sqrt{3}+\sqrt{2}}{\sqrt{3}-\sqrt{2}}=\ ?$} کے برابر کیا ہے::}{$5+2\sqrt{6}$}{1}{$\sqrt{6}$}{$5-2\sqrt{6}$}
\umcq{30}{اگر \lr{$x=1+\sqrt{2}$} ہو تو \lr{x$^{2}$=?}:}{$2\sqrt{2}$}{$1+2\sqrt{2}$}{3}{$3+2\sqrt{2}$}
\umcq{31}{اگر \lr{a+b=10،ab=21} ہو \lr{a$^{2}$+b$^{2}$=}:}{58}{100}{42}{121}
\umcq{32}{\lr{2$^{x}$=16} کا حل کیا ہے:}{6}{8}{2}{4}
\umcq{33}{\lr{16} کا جذر المربع کیا ہے:}{4}{-4}{\urdu{الف اور ب دونوں}}{\urdu{ان میں سے کوئی نہیں}}
\umcq{34}{\lr{$\sqrt{25}$} کی قیمت کیا ہے::}{5}{-5}{\urdu{الف اور ب دونوں}}{\urdu{ان میں سے کوئی نہیں}}
\umcq{35}{\lr{$F=\dfrac{9}{5}C+32$} کا استعمال کرتے ہوئے \lr{25} ڈگری سینٹی گریڈ کو فارن ہیٹ میں تبدیلی کریں:}{77$^{\circ}$F}{68$^{\circ}$F}{90$^{\circ}$F}{45$^{\circ}$F}
\umcq{36}{\lr{500} روپے کی لاگت پر \lr{20}٪ منافع کا مطلب ہے کہ اس کی قیمت فروخت ہے:}{\urdu{600روپے}}{\urdu{400روپے}}{\urdu{550روپے}}{\urdu{700روپے}}
\umcq{37}{\lr{2000} روپے پر \lr{2} سال کے لیے \lr{10}٪ کمپاؤنڈ سود ہے:}{\urdu{420روپے}}{\urdu{200روپے}}{\urdu{400روپے}}{\urdu{240روپے}}
\umcq{38}{اگر تین اعداد کا مجموعہ \lr{51} ہو تو،سب سے بڑا نمبر ہوگا:}{15}{16}{17}{18}
\umcq{39}{ایک مستطیل کی لمبائیاں \lr{$\sqrt{3}$} اور \lr{$\sqrt{12}$} ہو تو اس کا ایریا کیا ہوگا:: \lr{has an area of:5.svg"/>}}{$\sqrt{15}$}{6}{$3\sqrt{3}$}{12}
\umcq{40}{اگر \lr{5} پین کی قیمت \lr{45} روپے ہو تو، \lr{8} پن کی قیمت کیا ہوگی:}{100}{60}{90}{72}
\umcq{41}{ایک نمبر کو \lr{20} فیصد بڑھانے سے \lr{60} ہو جاتا ہے، اصل نمبر کیا ہوگا:}{50}{48}{60}{72}
\umcq{42}{دو نمبر \lr{3:4} کے تناسب میں ہیں۔ ان کا مجموعہ \lr{35} ہے۔ نمبر کیا ہوں گے:}{15,20}{12,16}{9,12}{18,24}
\umcq{43}{ایک ٹرین \lr{5} گھنٹوں میں \lr{300} کلو میٹر سفر کرتی ہے اس کی سپیڈ کیا ہوگی:}{50 km/h}{60 km/h}{45 km/h}{65 km/h}
\umcq{44}{سائیڈ \lr{$\sqrt{5}$} والے مربع کا رقبہ ہے::}{5}{10}{25}{20}
\umcq{45}{\lr{$\sqrt{0.25}+\sqrt{0.04}$} کا مجموعہ ہے::}{0.7}{0.5}{0.9}{1.0}
\umcq{46}{درجہ حرارت \lr{100} ڈگری سینٹی گریڈ کتنے کیلون کے برابر ہے:}{K 273}{K 373}{K 473}{K 173}
\umcq{47}{\lr{$\sqrt{7}$} ہے:}{\urdu{صحیح عدد}}{\urdu{ناطق عدد}}{\urdu{غیر ناطق عدد}}{\urdu{قدرتی عدد}}
\umcq{48}{\lr{$\pi$} اور \lr{e} ہیں:}{\urdu{قدرتی اعداد}}{\urdu{صحیح اعداد}}{\urdu{ناطق اعداد}}{\urdu{غیر ناطق اعداد}}
\umcq{49}{اگر \lr{n} ایک مکمل مربع نہیں ہے، تو \lr{$\sqrt{n}$} ہے::}{}{}{}{}
\umcq{50}{\lr{$\sqrt{3}+\sqrt{5}$} ہے:}{\urdu{مکمل عدد}}{\urdu{صحیح عدد}}{\urdu{ناطق عدد}}{\urdu{غیر ناطق عدد}}
\umcq{51}{تمام \lr{$x\in R$, $x=x$} کے لیے خاصیت کہلاتی ہے::}{}{}{}{}
\umcq{52}{فرض کریں \lr{$a,b,c\in R$} اگر \lr{$a>b$} اور \lr{$b>c\Rightarrow a>c$} یہ خاصیت ...... کہلاتی ہے:}{\urdu{ثلاثی}}{\urdu{متعدیت}}{\urdu{جمعی}}{\urdu{ضربی}}
\umcq{53}{\lr{$2^{x}\times 8^{x}=64$} ہو تو \lr{x} برابر ہے:}{$\dfrac{3}{2}$}{$\dfrac{3}{4}$}{$\dfrac{5}{6}$}{$\dfrac{2}{3}$}
\umcq{54}{فرض کریں \lr{$a,b\in R$} تو \lr{$a=b$} اور \lr{$b=a$} خاصیت ...... کہلاتی ہے:}{\urdu{عکسی}}{\urdu{تشاکل}}{\urdu{متعدیت}}{\urdu{جمعی}}
\umcq{55}{\lr{$\sqrt{75}+\sqrt{27}=$ ......}}{$\sqrt{102}$}{$9\sqrt{3}$}{$5\sqrt{3}$}{$8\sqrt{3}$}
\umcq{56}{\lr{$(3+\sqrt{5})(3-\sqrt{5})$} کا حاصل ضرب ہے::}{}{}{}{}
\umcq{57}{\lr{45000} کو سائنسی ترقیم میں تبدیل کریں:}{4.5 x 10$^{3}$}{4.5 x 10$^{4}$}{45 x 10$^{3}$}{0.45 x 10$^{5}$}

\umcq{58}{\lr{0.000078} کی سائنسی ترقیم ہے:}{$7.8\times 10^{-5}$}{$78\times 10^{-6}$}{$0.78\times 10^{-4}$}{$7.8\times 10^{5}$}
\umcq{59}{درج ذیل میں کون سا سائنسی ترقیم میں نہیں ہے:}{$3.4\times 10^{6}$}{$0.5\times 10^{3}$}{$6.02\times 10^{23}$}{$1.2\times 10^{-4}$}
\umcq{60}{\lr{$9.01\times 10^{5}$} کو عام ترقیم میں لکھیں::}{9010}{90100}{901000}{901}
\umcq{61}{\lr{$6.022\times 10^{23}$} میں عدد ہے:}{6}{6.022}{$10^{23}$}{23}
\umcq{62}{\lr{$73\times 10^{3}$} سائنسی ترقیم میں ہے:}{$7.3\times 10^{2}$}{$7.3\times 10^{3}$}{$7.3\times 10^{4}$}{$73\times 10^{3}$}
\umcq{63}{\lr{$6.7779\times 10^{3}$} کو سادہ شکل میں تبدیل کریں:}{677.76 km}{67.779 km}{67779 km}{6777.9 km}
\umcq{64}{\lr{$\log_{5}25=2$} کی قوت نمائی شکل ہے::}{25$^{2}$=5}{2$^{5}$=25}{5$^{2}$=25}{25$^{5}$=2}
\umcq{65}{\lr{10$^{3}$=1000} کی لوگارتھمی شکل ہے:}{$\log_{10}3=1000$}{$\log_{3}10=1000$}{$\log_{10}1000=3$}{$\log_{3}1000=10$}
\umcq{66}{حل کریں: \lr{$\log_{3}27=x$}:}{2}{9}{6}{3}
\umcq{67}{\lr{7$^{0}$=1} کو لوگارتھمی شکل میں تبدیل کریں:}{$\log_{1}7=0$}{$\log_{0}7=1$}{$\log_{7}1=0$}{$\log_{7}0=1$}
\umcq{68}{\lr{$\log_{10}100=2$} کی اساس ہے:}{2}{10}{100}{1}
\umcq{69}{\lr{log 0.0078} کی خاصیت ہوگی:}{-2}{-3}{3}{2}
\umcq{70}{\lr{log 345} کا منٹیسیا ہے: \lr{(given log 345=2.5378)}}{2}{0.5378}{5378}{2.5378}
\umcq{71}{اگر \lr{$\log x=\overline{3}.678$} ہو تو خاصیت ہوگی:}{3}{-3}{-4}{4}
\umcq{72}{اگر \lr{log x =2.3010}, ہو تو \lr{x=}:}{20}{200}{0.2}{2000}
\umcq{73}{قدرتی لوگارتھم کی اساس ہے:}{e}{10}{2}{3}
\umcq{74}{حل کریں: \lr{$\log(a)+\log(b)-\log(c)$}}{$\log\!\left(\dfrac{ab}{c}\right)$}{$\log(a+b-c)$}{$\log\!\left(\dfrac{ac}{b}\right)$}{$\log\!\left(\dfrac{a+b}{b}\right)$}
\umcq{75}{\lr{$\log(x^{2}y^{3})$} کو پھیلا کر لکھیں:}{6log}{logx + logy}{5log}{2logx + 3logy}
\umcq{76}{\lr{$\log_{3}\!\left(\dfrac{27}{9}\right)=$ ......}}{$\log_{3}27-\log_{3}9$}{$\log_{3}27+\log_{3}9$}{$\log_{3}(27\times 9)$}{$\log_{3}(27+3)$}
\umcq{77}{حل کریں: \lr{$\log_{5}2^{4}$}}{$4\log_{5}2$}{$2\log_{5}4$}{$\log_{5}8$}{$\log_{5}16$}
\umcq{78}{\lr{$\log_{4}8+\log_{4}2=$ : :}}{$\log_{8}2$}{$\log_{4}10$}{$\log_{8}4$}{$\log_{4}16$}
\umcq{79}{\lr{$\log_{7}5$} کو اساس \lr{10} میں تبدیل کریں::}{$\log_{7}10$}{$\dfrac{\log_{5}7}{\log_{5}10}$}{$\dfrac{\log_{10}5}{\log_{10}7}$}{$\log_{10}7$}
\umcq{80}{حل کریں: \lr{$\log_{6}6^{5}$}}{5}{6}{1}{0}
\umcq{81}{\lr{$\log_{3}81-\log_{3}9=$}}{$\log_{3}9$}{$\log_{3}72$}{$\log_{3}(81/9)$}{$\log_{3}(81\times 9)$}
\umcq{82}{اگر \lr{$\log_{a}b=c$} ہو تو:}{$c^{a}=b$}{$b^{c}=a$}{$a^{b}=c$}{$a^{c}=b$}
\umcq{83}{\lr{Ln(x)=2} کو قوت نمائی شکل میں تبدیل کریں:}{e$^{2}$=x}{x=10$^{2}$}{x=2$^{10}$}{x$^{2}$=e}
\umcq{84}{\lr{$M=\log(A/A_{o})$} کو استعمال کریں،اگر \lr{A=1000} اور \lr{A$_{0}$=10} ہو تو، \lr{M} معلوم کریں۔}{4}{3}{2}{5}
\umcq{85}{\lr{5.2 x 106} کی معیاری شکل ہے:}{52,000}{520,000}{5,200,000}{52,000,000}
\umcq{86}{\lr{0.00034} کی سائنسی ترقیم ہے:}{3.4 x 10$^{3}$}{3.4 x 10$^{-4}$}{3.4 x 10$^{4}$}{3.4 x 10$^{-3}$}
\umcq{87}{عام لوگارتھم کی اساس ہوتی ہے:}{2}{10}{5}{e}

\umcq{88}{\lr{Log$_{2}$2$^{3}$}:}{1}{2}{5}{3}
\umcq{89}{\lr{Log 100 = ......}:}{2}{3}{10}{1}
\umcq{90}{اگر \lr{log 2 = 0.3010} ہو تو \lr{log 200} ہے:}{1.3010}{0.6010}{2.3010}{2.6010}
\umcq{91}{\lr{Log (0) = ......}:}{Positive}{Negative}{Zero}{Undefined}
\umcq{92}{\lr{log 10,000 = .......}:}{2}{3}{4}{5}
\umcq{93}{\lr{log 5 + log 3 =......}:}{Log 0}{Log 2}{$\log\!\left(\dfrac{5}{3}\right)$}{Log 15}
\umcq{94}{\lr{3$^{4}$= 81} کی لوگارتھمی شکل ہے:}{log$_{3}$4=81}{log$_{4}$3 = 81}{log$_{3}$81 = 4}{log$_{4}$81 = 3}
\umcq{95}{سلسلہ \lr{0،1،1،2،3،5،8،13،21 .......} ...... سلسلہ کی مثال ہے:}{\urdu{رومن}}{\urdu{مفرد}}{\urdu{طاق}}{\urdu{فیبوناکی}}
\umcq{96}{جارج کینٹر ...... ریاضی دان تھا:}{\urdu{جاپانی}}{\urdu{جرمن}}{\urdu{یونانی}}{\urdu{مصری}}
\umcq{97}{فیبونا کی سلسلہ کا فارمولا \lr{F$_{n}$=......}:}{F$_{n-1}$+F$_{n+2}$}{F$_{n-1}$+F$_{n-2}$}{F$_{n-1}$-F$_{n+2}$}{F$_{n-1}$+F$_{n}$}
\umcq{98}{سیٹ کے نام ظاہر کرتے ہیں:}{1,2,3,..}{A,B,C ...}{\urdu{الف اور ب دونوں}}{\urdu{ان میں سے کوئی نہیں}}
\umcq{99}{ایک سیٹ کو بیان کرنے کے کتنے طریقے ہیں:}{1}{2}{3}{4}
\umcq{100}{سیٹ \lr{A=\{a,e,i,o,u\}} کس انداز میں لکھا گیا ہے:}{\urdu{بیانیہ طریقہ}}{\urdu{اندراجی طریقہ}}{\urdu{ترقیم سیٹ ساز}}{\urdu{ان میں سے کوئی نہیں}}
\umcq{101}{سیٹ سے تعلق کے لیے کون سا سمبل استعمال کیا جاتا ہے:}{$\subseteq$}{$\in$}{$\notin$}{$\forall$}
\umcq{102}{سیٹ \lr{\{2,3,5,7,11,......\}} ...... رکن کہلاتا ہے:}{\urdu{قدرتی}}{\urdu{مکمل}}{\urdu{طاق}}{\urdu{پرائم}}
\umcq{103}{\lr{Q$\cup$Q$'$=......}:}{N}{R}{Q}{Q$'$}
\umcq{104}{سیٹ \lr{\{0\}} ہے:}{\urdu{یک رکنی}}{\urdu{خالی}}{\urdu{صفر}}{\urdu{پرائم}}
\umcq{105}{ایسا سیٹ جس میں کوئی رکن نہ ہو...... سیٹ کہلاتا ہے:}{\urdu{یک رکنی}}{\urdu{خالی}}{\urdu{پاور}}{\urdu{پرائم}}
\umcq{106}{خالی سیٹ کے لیے سمبل استعمال کیا جاتا ہے:}{$\{\phi\}$}{$\phi$}{\{\{ \}\}}{None of these}
\umcq{107}{مساوی سیٹ کی مثال ہے:}{\{a,b,c\}, \{1,2,3\}}{\{1,2,3\}, \{2,1,3\}}{\urdu{الف اور ب دونوں}}{\urdu{ان میں سے کوئی نہیں}}
\umcq{108}{مترادف سیٹ \lr{A} اور \lr{B} کے لیے سمبل استعمال کیا جاتا ہے:}{A=B}{$A\equiv B$}{A+B}{AUB}
\umcq{109}{اگر \lr{A} سیٹ \lr{B} کا سب سیٹ ہو تو اس کو کس طرح لکھا جائے گا:}{$A\subseteq B$}{$A\supseteq B$}{$A\in B$}{$A\notin B$}
\umcq{110}{پاور سیٹ \lr{C=\{a,b,c,d\}} کے ارکان ہیں:}{8}{16}{24}{26}
\umcq{111}{اگر \lr{D=\{a\}} ہو تو \lr{P(D)=......}:}{\{a\}}{$\phi,\{a\}$}{$\{\phi,\{a\}\}$}{\urdu{ان میں سے کوئی نہیں}}
\umcq{112}{\lr{S$_{1}$=} طاق اعداد کا سیٹ،\lr{S$_{2}$=} جفت اعداد کا سیٹ ہے تو دونوں کا تقاطع کیا ہوگا:}{\urdu{یونیورسل سیٹ}}{\urdu{خالی سیٹ}}{\urdu{غیر مشترک}}{\urdu{مفرد سیٹ}}
\umcq{113}{\lr{$\{x\mid x\in A\ \vee\ x\in B\}$} کو اس طرح ظاہر کیا جا سکتا ہے:}{AUB}{$A\cap B$}{A-B}{$\phi$}
\umcq{114}{تقاطع کے لیے علامت استعمال کی جاتی ہے:}{U}{I}{$\cap$}{-}
\umcq{115}{\lr{$A\cap B=......$}}{$\{x\mid x\in A\vee x\in B\}$}{$\{x\mid x\in A\wedge x\in B\}$}{$\{x\mid x\in A\wedge x\notin B\}$}{$\{x\mid x\notin A\wedge x\in B\}$}
\umcq{116}{\lr{$A\cap B$=\{1,2,3\}} اگر \lr{A=\{1,2,3,7\}} ہو تو \lr{B=......}}{\{1,2,6\}}{\{1,2,3,4\}}{\{1,3,6\}}{\{1,6,7\}}
\umcq{117}{اگر \lr{$A\cap B=\phi$} ہو تو \lr{A} اور \lr{B} سیٹ ہونگے::}{\urdu{سب سیٹ}}{\urdu{متراکب سیٹ}}{\urdu{غیر مشترک}}{\urdu{ان میں سے کوئی نہیں}}
\umcq{118}{متراکب سیٹ کی مثال کونسی ہے:}{\{1,2\}, \{2,3,4\}}{\{1,2\}, \{1,2,3,4\}}{\{1,2\}, \{3,4\}}{None of these}
\umcq{119}{\lr{$\{x\mid x\in A\vee x\notin B\}$} کو ظاہر کیا جاتا ہے:}{A-B}{B-A}{AUB}{A+B}
\umcq{120}{اگر \lr{A=\{1,2,3,4,5\}} اور \lr{B=\{4,5,6,7,8,9,10\}} ہو تو \lr{A-B} ہو گا:}{\{4,5\}}{\{1,2,3\}}{\{7,8,9,10\}}{\{4,5,6,7,8,9,10\}}
\umcq{121}{\lr{A$'$......}:}{U-A}{A-U}{U}{U+A}

\umcq{122}{\lr{A$'$=......}:}{U-A}{A-U}{U}{$\phi$}
\umcq{123}{\lr{(AUB)$'$=......}:}{$A\cap B$}{A$'$UB$'$}{$A'\cap B'$}{A$'$UB}
\umcq{124}{\lr{$A\cap B=B\cap A$} تقاطع کی ...... خاصیت کہلاتی ہے::}{\urdu{خاصیت تلازم}}{\urdu{خاصیت مبادلہ}}{\urdu{خاصیت تقسیمی}}{\urdu{کوئی نہیں}}
\umcq{125}{اگر \lr{A=\{1,3,5,7,9,11\}} ہو تو، \lr{n(A)=.......}:}{6}{4}{5}{2}
\umcq{126}{\lr{$n(A)+n(B)-n(A\cap B)=......$}}{$n(A\cap B)$}{$n(A\cup B)$}{n}{None of these}
\umcq{127}{\lr{A x B} کو علامتی طور پر لکھا جاتا ہے:}{$\{(x,y)\mid x\in A$ and $x\in B\}$}{$\{(x,y)\mid x\in A$ and $y\in B\}$}{$\{(x,y)\mid x\notin A$ and $x\in B\}$}{$\{(x,y)\mid x\in B$ and $y$}
\umcq{128}{فرض کریں \lr{A=\{1,2,3\}} ہو تو خاصیت معلوم کریں: \lr{$r=\{(x,y\mid x<y)\}$}}{\{(2,1),(3,1),(3,2)\}}{\{(1,2),(1,3),(2,3)\}}{\{(1,1),(2,2),(3,3)\}}{None of these}
\umcq{129}{اگر \lr{r=\{(1,2),(1,3),(2,3)\}} ہو تو ڈومین \lr{r=}:}{\{1,2\}}{\{1\}}{\{2,3\}}{\{1,2,3\}}
\umcq{130}{ہر مترتب جوڑا...... کو آرڈینٹ پر مشتمل ہوتا ہے:}{1}{2}{3}{4}
\umcq{131}{مترتب جوڑے \lr{(x,y)} میں \lr{x} کہلاتا ہے:}{\urdu{ابسیسہ}}{\urdu{آرڈینیٹ}}{\urdu{دونوں}}{\urdu{کوئی نہیں}}
\umcq{132}{اگر ایک تفاعل \lr{$f:A\rightarrow B$} اس طرح ہو کہ ریج \lr{$f\neq B$}, تو \lr{f} ...... تفاعل کہلائے گا:}{\urdu{آن ٹو}}{\urdu{ان ٹو}}{\urdu{ون-ون}}{\urdu{اجیکٹیو}}
\umcq{133}{اگر \lr{$f(n)=2n-1$} ہو تو \lr{$f(7)$} ہے:}{12}{13}{14}{15}
\umcq{134}{اگر \lr{$g(n)=x^{2}-3$} ہو تو \lr{g(4)} ہوگا:}{7}{-7}{13}{9}
\umcq{135}{سیٹ \lr{$\left\{1,\dfrac{1}{3},\dfrac{1}{5},\dfrac{1}{7},\ldots\right\}$} کی ترقیم ساز صورت ہے::}{$\left\{x\mid x=\dfrac{1}{n},n\in W\right\}$}{$\left\{x\mid x=\dfrac{1}{2n+1},n\in W\right\}$}{$\left\{x\mid x=\dfrac{1}{n+1},n\in W\right\}$}{$\{x\mid x=2n+1,\ n\in$}
\umcq{136}{اگر \lr{A = \{ \}} ہو تو \lr{P(A)} برابر ہے:}{\{ \}}{\{1\}}{\{\{ \}\}}{$\phi$}
\umcq{137}{اگر \lr{$U=\{1,2,3,4,5\}$}، \lr{$A=\{1,2,3\}$} اور \lr{$B=\{3,4,5\}$} ہو تو \lr{$U-(A\cap B)$} برابر ہے:}{\{1,2,4,5\}}{\{2,3\}}{\{1,3,4,5\}}{\{1,2,3\}}
\umcq{138}{اگر \lr{A} اور \lr{B} متراکب سیٹ ہوں تو \lr{$n(A-B)$} برابر ہے::}{n(A)}{n(B)}{$A\cap B$}{$n(A)-n(A\cap B)$}
\umcq{139}{اگر \lr{$A\subseteq B$} اور \lr{$B-A\neq\phi$} ہو تو \lr{$n(B-A)$} برابر ہے:}{0}{-n(B)}{n(A)}{n(B)-n(A)}
\umcq{140}{اگر \lr{$n(A\cup B)=50,n(A)=30$} اور \lr{$n(B)=35$} ہو تو \lr{$n(A\cap B)$} برابر ہے:}{23}{15}{9}{40}
\umcq{141}{اگر \lr{$A=\{1,2,3,4\}$} اور \lr{$B=\{x,y,z\}$} ہو تو سیٹ \lr{A} اور \lr{B} کے کارتیسی حاصل ضرب میں کل ..... ارکان ہیں؟}{13}{12}{10}{6}
\umcq{142}{\lr{$f(x)=x^{2}-3x+2$} ہو تو \lr{$f(a+1)$} کی قیمت ہے:}{a+1}{a$^{2}$+1}{a$^{2}$+2a+1}{a$^{2}$-a}
\umcq{143}{فرض کریں \lr{$f(x)=3x+1$} اگر \lr{$f(x)=28$} ہو تو \lr{x} کی قیمت معلوم کریں:}{9}{27}{3}{18}
\umcq{144}{اگر \lr{$A=\{1,2,3\}$} اور \lr{$B=\{a,b\}$} دو غیر خالی سیٹ ہیں اور \lr{$f:A\rightarrow B$} ایک تفاعل ہے جس کو اس طرح بیان کیا گیا ہو۔ \lr{$f=\{(1,a),(2,b),(3,b)\}$} تو درج ذیل میں سے کون سا انتخاب درست ہے؟}{\urdu{f ان جیکٹو ہے}}{\urdu{f سر جیکٹو ہے}}{\urdu{f بائی جیکٹو ہے}}{\urdu{f صرف ان ٹو ہے}}
\umcq{145}{\lr{2x - 6} کا مشترک جزو ضربی ہے:}{2}{3}{x}{None of these}
\umcq{146}{\lr{x$^{2}$-3x-10} کا جزو ضربی ہے:}{(x + 2)(x + 5)}{(x + 2)(x - 5)}{(x - 2)(x + 5)}{(x - 2)(x - 5)}
\umcq{147}{کونسی مساوات دو درجی مساوات کی قسم ہے:}{ax + b}{ax$^{2}$+ bx + c}{ax$^{3}$+ bx$^{2}$+ c}{None of these}
\umcq{148}{\lr{x$^{2}$-11x +24} کی تجزی کے لیے کونسا جوڑا درست ہے:}{+3, +8}{-3,-8}{-3, +8}{+3,-8}
\umcq{149}{دو درجے رکھنے والی مساوات ...... مساوات کہلاتی ہے:}{\urdu{کیوبک}}{\urdu{لینیئر}}{\urdu{دو درجی}}{\urdu{ان میں سے کوئی نہیں}}
\umcq{150}{\lr{3x$^{2}$-4x-4} کے جزو ضربی ہیں:}{(3x + 2)(x - 2)}{(3x - 2)(x + 2)}{(3x + 2)(x + 2)}{None of these}
\umcq{151}{درج ذیل میں کونسا پرائم نمبر ہے:}{9}{15}{13}{21}

\umcq{152}{\lr{x$^{2}$-2ab+b$^{2}$} برابر ہے:}{(a + b)$^{2}$}{(a - b)$^{2}$}{a$^{2}$- b$^{2}$}{a$^{2}$+ b$^{2}$}
\umcq{153}{\lr{8$^{3}$+27} کا جزو ضربی ہے:}{(2x + 3)(4x$^{2}$+ 6x + 9)}{(2x + 3)(4x$^{2}$- 6x + 9)}{(2x - 3)(4x$^{2}$+ 6x + 9)}{(2x - 3)(4x$^{2}$- 6x + 9)}
\umcq{154}{\lr{2x$^{2}$+7x+3} میں ایک جزو ضربی \lr{2x + 1} ہو تو دوسرا کیا ہوگا:}{(x + 2)}{(x + 1)}{(x + 3)}{(x + 4)}
\umcq{155}{ایسی مساوات جس کی دو رقمیں ہوں کہلاتی ہیں:}{\urdu{یک رقمی}}{\urdu{دو رقمی}}{\urdu{تین رقمی}}{\urdu{ان میں سے کوئی نہیں}}
\umcq{156}{\lr{a$^{3}$+b$^{3}$+3ab(a+b)} برابر ہے:}{(a + b)$^{3}$}{(a - b)$^{3}$}{a$^{3}$- b$^{3}$}{a$^{3}$+ b$^{3}$}
\umcq{157}{عادا اعظم کو حل کرنے کے کتنے طریقے ہیں:}{1}{2}{3}{4}
\umcq{158}{\lr{6x$^{2}$y,9xy$^{2}$} کا عادا اعظم ہے:}{3x$^{2}$y}{3xy}{6x$^{2}$y}{54x$^{2}$y$^{3}$}
\umcq{159}{ذواضعاف اقل \lr{=......}:}{$\dfrac{p(x)\times q(x)}{H.C.F}$}{$\dfrac{H.C.F\ \times\ p(x)}{q(x)}$}{P}{None of these}
\umcq{160}{\lr{a$^{2}$-2a+1} کا جذر المربع ہے:}{$\pm$(a-1)}{$\pm$(a+1)}{a - 1}{a + 1}
\umcq{161}{\lr{x$^{2}$-8x+16} کا جذر المربع ہے:}{(x+4)}{-(x+4)}{\urdu{A اور B دونوں}}{\urdu{ان میں سے کوئی بھی نہیں}}
\umcq{162}{انویسٹمنٹ پروا پسی \lr{R(x)} ہے \lr{R(x) =x$^{2}$+6x-8}۔ انویسٹمنٹ لیولز کیا ہیں جن میں نتیجہ صفر ہے:}{x=2, x=4}{x=1, x=8}{x=2, x=-4}{x=-2, x=4}
\umcq{163}{\lr{12x +36} کی تجزی ہے:}{12(x+3)}{12(3x)}{12(3x+1)}{x(12+36x)}
\umcq{164}{\lr{$4x^{2}-12y+9$} کے اجزائے ضربی ہیں::}{$(2x+3)^{2}$}{$(2x-3)^{2}$}{(2x-3) (2x+3)}{$(2+3x)(2-3x)^{2}$}
\umcq{165}{\lr{$a^{3}b^{3}$} اور \lr{$ab^{2}$} کا عادا اعظم ہے:}{$a^{3}b^{3}$}{$ab^{2}$}{$a^{4}b^{5}$}{$a^{2}b$}
\umcq{166}{\lr{$16x^{2}$}،\lr{4x} اور \lr{30xy} کا ذواضعاف اقل ہے::}{$480x^{3}y$}{240xy}{$240x^{2}y$}{$120x^{4}y$}
\umcq{167}{دو کثیر رقمی جملوں کا\lr{......}= ذواضعاف اقل اور عاد اعظم کا حاصل ضرب:}{\urdu{مجموعہ}}{\urdu{فرق}}{\urdu{حاصل ضرب}}{\urdu{حاصل تقسیم}}
\umcq{168}{\lr{$x^{2}-6x+9$} کا جذر المربع ہے::}{$\pm(x-3)$}{$\pm(x+3)$}{x-3}{x+3}
\umcq{169}{\lr{$(a-b)^{2}$} اور \lr{$(a-b)^{4}$} کا ذواضعاف اقل ہے:}{$(a-b)^{2}$}{$(a-b)^{3}$}{$(a-b)^{4}$}{$(a-b)^{6}$}
\umcq{170}{\lr{$x^{3}+3x^{2}+3x+1$} کی تجزی ہے::}{$(x+1)^{3}$}{$(x-1)^{3}$}{$(x+1)(x^{2}+x+1)$}{$(x-1)(x^{2}-x+1)$}
\umcq{171}{سہ درجی کثیر رقمی جملے کا درجہ ہوتا ہے:}{1}{2}{3}{4}
\umcq{172}{\lr{$x^{3}-27$} کا ایک جزو ضربی ہے::}{x-3}{x+3}{$x^{2}-3x+9$}{Both a and c}
\umcq{173}{درج ذیل میں سے کونسی مساوات ایک متغیر میں یک درجی مساوات ہے:}{2x + 3y = 7}{5x - 4 = 0}{x$^{2}$+ 3x = 8}{xy + 2 = 5}
\umcq{174}{مساوات \lr{3x - 7 = 2} کا حل ہے:}{3}{-3}{$\dfrac{5}{3}$}{9}
\umcq{175}{غیر درجی مساوات \lr{4x+2>10}, میں \lr{x} کا حل کیا ہے:}{x $>$ 2}{x $<$ 2}{x $>$1}{x $<$ 2}
\umcq{176}{\lr{:}}{$<$}{$>$}{$\leq$}{$\geq$}
\umcq{177}{دو متغیرات میں یک درجی مساوات کو گراف کی شکل میں اس طرح لکھا جاتا ہے:}{\urdu{ترچھا}}{\urdu{سیدھی لائن}}{\urdu{ایک پیرابولا}}{\urdu{دائرے میں}}
\umcq{178}{\lr{2(x-3)=4} کا حل کیا ہے:}{x = 2}{x = 5}{x = 3}{x = 1}
\umcq{179}{غیر مساوات \lr{$x-5\geq 3$} کا مطلب ہے::}{\lr{x 3} \urdu{سے بڑا ہے}}{\lr{x 3} \urdu{سے بڑا اور برابر بھی ہے}}{\lr{x, 8} \urdu{سے بڑا اور برابر ہے}}{\lr{x, 3} \urdu{سے چھوٹا ہے}}
\umcq{180}{مساوات \lr{x+2y=6} ظاہر کرتی ہے:}{\urdu{دو درجی مساوات}}{\urdu{یک درجی مساوات}}{\urdu{کفایتی مساوات}}{\urdu{ان میں سے کوئی نہیں}}

\umcq{181}{غیر مساوات \lr{$x+2y\leq 8$} کی متعلقہ مساوات ہے::}{x + 2y = 8}{x + 2y $>$ 8}{x + 2y $<$ 8}{$x+2y\geq 8$}
\umcq{182}{سسٹم \lr{x-y=2} اور \lr{x+y=6} کا حل ہے:}{(4 ,2)}{(2 ,4)}{(3 ,3)}{(6 ,0)}
\umcq{183}{اگر \lr{x=3, 5x-2=k} کا حل ہے تو \lr{k} ہو گا:}{11}{13}{15}{17}
\umcq{184}{یک درجی غیر مساوات کے سسٹم کی گرافیکل نمائندگی کیا ہوگی:}{\urdu{سایہ دار علاقہ}}{\urdu{ایک نقطہ}}{\urdu{ایک سیدھی لائن}}{\urdu{ایک کرو}}
\umcq{185}{مساوات \lr{2x-y=4} کو کونسی غیر مساوات کے برابر ہے:}{2x - y $>$ 4}{2x - y $<=$ 4}{2x - y $>=$ 4}{None of these}
\umcq{186}{درج ذیل میں سے کونسا آبجیکٹو تفاعل ہے:}{x + y}{x$^{2}$ + y$^{2}$}{3x + 2y}{Both a and c}
\umcq{187}{یک درجی غیر مساوات کے نظام میں قابل عمل خطہ ہمیشہ ہوتا ہے:}{\urdu{ایک نقطہ}}{\urdu{ایک پلین کا علاقہ}}{\urdu{ایک خالی سیٹ}}{\urdu{ایک لائن کا سلسلہ}}
\umcq{188}{ایک عمودی لائن پلین کو تقسیم کرتی ہے:}{\urdu{دو آدھے پلین میں}}{\urdu{چار حصوں میں}}{\urdu{تین سیکشنز میں}}{\urdu{لامحدود حصے}}
\umcq{189}{اگر کسی غیر مساوات کے نظام کا کوئی حل نہ ہو تو اس کا قابل عمل خطہ ہو گا:}{\urdu{لامتناہی}}{\urdu{ایک نقطہ}}{\urdu{ایک سنگلحصہ}}{\urdu{خالی}}
\umcq{190}{دو \lr{Half-planes} کا امتزاج "خط" کہلاتا ہے:}{A vertex}{A feasible region}{A line}{A parabola}
\umcq{191}{\lr{$x+y\geq 5$} کے گراف میں لائنیں شامل ہوتی ہیں:}{x + y = 5}{x + y $<$ 5}{x + y $>$ 5}{None of these}
\umcq{192}{یک درجی مسئلے میں نہایت موافق حل معلوم کرنے کے لیے کونسا طریقہ استعمال کیا جاتا ہے:}{\urdu{گرافیکل طریقہ}}{\urdu{الجبری طریقہ}}{\urdu{سادہ ترین طریقہ}}{\urdu{یہ تمام}}
\umcq{193}{مساوات \lr{x=2} ظاہر کرتی ہے:}{\urdu{افقی لکیر}}{\urdu{عمودی لکیر}}{\urdu{ترچھی لکیر}}{\urdu{ان میں سے کوئی نہیں}}
\umcq{194}{ایک پوائنٹ جو دی ہوئی غیر مساواتوں کو مطمئین کرتا ہے، کہلاتا ہے:}{\urdu{قابل عمل خطہ}}{\urdu{نہایت موافق حل}}{\urdu{باونڈری پوائنٹ}}{\urdu{ان میں سے کوئی نہیں}}
\umcq{195}{لائن \lr{y=3x+2} سلوپ رکھتی ہے:}{2}{3}{-3}{-2}
\umcq{196}{اگر غیر مساواتیں \lr{$x\geq 0$} اور \lr{$y\geq 0$} دی گئی ہوں تو حل شدہ خطہ کس میں ہوگا؟}{\urdu{پہلے ربع میں}}{\urdu{دوسرے ربع میں}}{\urdu{تیسرے ربع میں}}{\urdu{چوتھے ربع میں}}
\umcq{197}{غیر مساوات \lr{5x - 3y $<$ 15} ظاہر کرتی ہے:}{\urdu{ایک پیرابولا}}{\urdu{ایک سیدھی لائن}}{\urdu{ایک سایہ دار علاقہ}}{\urdu{ایک سنگل پوائنٹ}}
\umcq{198}{ایک تفاعل کی زیادہ یا کم سے کم قیمت رونما ہوتی ہے:}{\urdu{رینڈم نقطے پر}}{\urdu{حل شدہ نقطے پر}}{\urdu{کونے والے نقطے پر}}{\urdu{نامعلوم نقطے پر}}
\umcq{199}{اگر \lr{$x-y\leq 4$} ہو تو حل کا علاقہ ہوگا:}{\urdu{لائن کے اوپر}}{\urdu{لائن کے نیچے}}{\urdu{لائن پر}}{\urdu{ان میں سے کوئی نہیں}}
\umcq{200}{\lr{x+y=4} اور \lr{x-y=2} کا قاطع ہوگا:}{(1,3)}{(3,1)}{(2,2)}{(4,0)}
\umcq{201}{درج ذیل میں یک درجی مساوات ہے:}{5x $>$ 7}{4x - 2 $<$ 1}{2x + 1 = 1}{4 = 1 + 3}
\umcq{202}{\lr{5x - 10 = 10} کا حل ہے:}{0}{50}{4}{-4}
\umcq{203}{اگر \lr{7x + 4 $<$ 6x + 6} ہو تو \lr{x} کا حل ہو گا؟:}{$(2,\infty)$}{$[2,\infty)$}{$(-\infty,2)$}{$(-\infty,2]$}
\umcq{204}{عمودی خط مستوی کو تقسیم کرتا ہے:}{\urdu{بائیں نصف مستوی}}{\urdu{دائیں نصف مستوی}}{\urdu{پوری مستوی}}{\urdu{دو آدھی مستویاں}}
\umcq{205}{یک درجی مساوات سے بنے والی مساوات کو کہا جاتا ہے:}{\urdu{سہ درجی مساوات}}{\urdu{متعلقہ مساوات}}{\urdu{دو درجی مساوات}}{\urdu{قابل عمل خطہ}}
\umcq{206}{\lr{3x + 4 $<$ 0} ہے:}{\urdu{مساوات}}{\urdu{غیر مساوات}}{\urdu{غیر مساوات نہیں}}{\urdu{اکائی}}
\umcq{207}{کونے کا نقطہ کہلاتا ہے:}{\urdu{کوڈ}}{\urdu{راس}}{\urdu{خط}}{\urdu{خطہ}}
\umcq{208}{\lr{(0,0)} کا حل ہے:}{4x + 5y $>$ 8}{3x + y $>$ 6}{-2x + 3y $<$ 0}{x + y $>$ 4}
\umcq{209}{حل کا نقطہ جو پہلے ربع تک محدود ہو کہلاتا ہے:}{\urdu{آبجیکٹو خطہ}}{\urdu{قابل عمل خطہ}}{\urdu{حل خطہ}}{\urdu{رکاوٹوں کا خطہ}}
\umcq{210}{تفاعل جس کو کم سے کم یا زیادہ سے زیادہ کیا جائے، کہلاتا ہے:}{\urdu{حل کا تفاعل}}{\urdu{آبجیکٹو تفاعل}}{\urdu{قابل عمل تفاعل}}{\urdu{ان میں سے کوئی نہیں}}
\umcq{211}{\lr{$\angle AOB$} میں عمود ہے:}{A}{O}{B}{None of these}
\umcq{212}{رے \lr{OA} کو علامتی طور پر لکھا جا سکتا ہے:}{$\overleftrightarrow{OA}$}{$\widehat{OA}$}{$\overline{OA}$}{$\overrightarrow{OA}$}
\umcq{213}{\lr{65$^{\circ}$} کا کو ٹرمینل ہے:}{360$^{\circ}$}{-360$^{\circ}$}{425$^{\circ}$}{-405$^{\circ}$}
\umcq{214}{ایک دائرے کا مکمل چکر.............ریڈیئن کے برابر ہوتا ہے۔:}{$\pi$}{$2\pi$}{$3\pi$}{$4\pi$}
\umcq{215}{سیکٹر کا ایریا ...... ہے:}{$\dfrac{1}{2}r^{2}\theta$}{$\pi r^{2}$}{$4\pi r^{2}$}{$\dfrac{1}{2}r^{3}\theta$}

\umcq{216}{\lr{$\dfrac{\pi}{4}$} ریڈین کی قدر ہے:}{$\sqrt{2}$}{$\sqrt{3}$}{$\dfrac{1}{\sqrt{2}}$}{$\dfrac{1}{\sqrt{3}}$}
\umcq{217}{ڈگری میں مکمل روٹیشن کی پیمائش کیا ہے:}{90$^{\circ}$}{180$^{\circ}$}{270$^{\circ}$}{360$^{\circ}$}
\umcq{218}{ریڈین میں قائمہ الزاویہ کی پیمائش کیا ہے:}{$\dfrac{\pi}{2}$}{$\pi$}{$\dfrac{3\pi}{2}$}{$2\pi$}
\umcq{219}{درج ذیل میں سے منفرجہ زاویہ ہے:}{45$^{\circ}$}{90$^{\circ}$}{120$^{\circ}$}{180$^{\circ}$}
\umcq{220}{ریڈین میں سیدھے زاویے کی پیمائش کیا ہے:}{$2\pi$}{$\dfrac{3\pi}{2}$}{$\pi$}{$\dfrac{\pi}{2}$}
\umcq{221}{ڈگری میں \lr{1} ریڈین کی پیمائش کیا ہے:}{57.2958$^{\circ}$}{90$^{\circ}$}{180$^{\circ}$}{360$^{\circ}$}
\umcq{222}{\lr{73.12$^{\circ}$} کو ڈگری، منٹ اور سیکنڈ میں تبدیل کریں:}{73$^{\circ}$12$'$ 12$''$}{73$^{\circ}$12$'$ 7$''$}{73$^{\circ}$7$'$ 7$''$}{73$^{\circ}$7$'$ 12$''$}
\umcq{223}{\lr{45$^{\circ}$ 45$'$ 45$''$} کو اعشاری ڈگری میں تبدیل کریں:}{}{}{}{}
\umcq{224}{\lr{1$^{\circ}$} کی سیکنڈ میں کیا قدر رہے گی:}{60$''$}{360$''$}{3600$''$}{60$'$}
\umcq{225}{پورے دائرے کی ریڈین میں کیا پیمائش ہے:}{$2\pi$}{$\dfrac{3\pi}{2}$}{$\pi$}{$\dfrac{\pi}{2}$}
\umcq{226}{\lr{$\dfrac{5\pi}{3}$} ریڈین کو ڈگری میں تبدیل کریں:}{300$^{\circ}$}{150$^{\circ}$}{210$^{\circ}$}{330$^{\circ}$}
\umcq{227}{\lr{15$^{\circ}$} کو ریڈین میں تبدیل کریں۔:}{$\dfrac{\pi}{12}$}{$\dfrac{\pi}{6}$}{$\dfrac{\pi}{4}$}{$\dfrac{\pi}{3}$}
\umcq{228}{\lr{1.2} ریڈین کو ڈگری میں تبدیل کریں:}{90$^{\circ}$}{57.2958$^{\circ}$}{45$^{\circ}$}{68.75$^{\circ}$}
\umcq{229}{قائمہ الزاویہ مثلث \lr{ABC} میں \lr{$m\angle B=90^{o}$} ہو تو کون سا بیان درست ہو گا:}{a$^{2}$+ b$^{2}$= c$^{2}$}{a$^{2}$- b$^{2}$= c$^{2}$}{a$^{2}$+ b$^{2}$= b$^{2}$}{a$^{2}$+ b$^{2}$=a$^{2}$}
\umcq{230}{\lr{sin$\theta$} کا ریسی پروکل ہے::}{$\sec\theta$}{$\mathrm{cosec}\,\theta$}{$\cot\theta$}{$\tan\theta$}
\umcq{231}{\lr{cot $\theta$ = ...................}.}{$\dfrac{1}{\sin\theta}$}{$\dfrac{1}{\cos\theta}$}{$\dfrac{1}{\tan\theta}$}{1-cos$^{2}\theta$}
\umcq{232}{\lr{sin $\theta$ = ..................}.}{$\sqrt{1-\cos^{2}\theta}$}{$\dfrac{1}{\tan\theta}$}{$\dfrac{\sin\theta}{\cos\theta}$}{$1-\cos^{2}\theta$}
\umcq{233}{اگر \lr{$\sin\theta=\dfrac{2}{3}$} ہو تو \lr{cosec $\theta$} کی قدر کیا ہوگی:}{$\dfrac{2}{3}$}{$\dfrac{-2}{3}$}{$\dfrac{3}{2}$}{$\dfrac{-3}{2}$}
\umcq{234}{بنیادی تکونیاتی نسبتیں ہیں:}{3}{4}{5}{6}
\umcq{235}{\lr{cot 30$^{\circ}$} کی قدر کیا ہوگی؟}{$\sqrt{3}$}{$\dfrac{1}{\sqrt{3}}$}{1}{0}
\umcq{236}{\lr{sin (90$^{\circ}$-$\theta$)} کی قدر کیا ہوگی:}{Tan$\theta$}{Sin$\theta$}{Cos$\theta$}{Cot$\theta$}
\umcq{237}{\lr{sin $\theta$} اور \lr{cos $\theta$} کی صورت میں \lr{tan $\theta$} کی قدر کیا ہوگی:}{$\dfrac{\sin\theta}{\cos\theta}$}{$\dfrac{\cos\theta}{\sin\theta}$}{$\dfrac{1}{\sin\theta}$}{$\dfrac{1}{\cos\theta}$}
\umcq{238}{درج ذیل میں سے کون پاستھا گورین اکائی ہے:}{$\sin^{2}\theta+\cos^{2}\theta=1$}{$\sin\theta+\cos\theta=1$}{$\tan\theta+\cot\theta=1$}{$\sec\theta+\mathrm{cosec}\,\theta=1$}
\umcq{239}{\lr{tan$^{2}$ $\theta$ + 1} کی قدر کیا ہوگی؟}{Sin$^{2}\theta$}{Cosec$^{2}\theta$}{Cot$^{2}\theta$}{Sec$^{2}\theta$}
\umcq{240}{\lr{1 + cot$^{2}\theta$} کی قدر کیا ہوگی؟}{Cosec$^{2}\theta$}{Sec$^{2}\theta$}{Tan$^{2}\theta$}{Sin$^{2}\theta$}
\umcq{241}{اگر \lr{$\tan\theta=\dfrac{3}{4}$} ہو تو، \lr{sin$\theta$} کی قدر کیا ہوگی؟}{$\dfrac{5}{3}$}{$\dfrac{4}{5}$}{$\dfrac{3}{5}$}{$\dfrac{5}{4}$}
\umcq{242}{\lr{sin 45$^{\circ}$} کی قدر کیا ہوگی:}{$\dfrac{1}{\sqrt{2}}$}{$\dfrac{\sqrt{3}}{2}$}{$\dfrac{1}{2}$}{1}

\umcq{243}{\lr{cos 30$^{\circ}$} کی قدر کیا ہوگی؟:}{$\dfrac{1}{2}$}{$\dfrac{\sqrt{3}}{2}$}{$\dfrac{1}{\sqrt{2}}$}{1}
\umcq{244}{\lr{tan 60$^{\circ}$} کی قدر کیا ہوگی:}{$\sqrt{3}$}{$\dfrac{1}{\sqrt{3}}$}{1}{$\dfrac{\sqrt{3}}{2}$}
\umcq{245}{\lr{cot 45$^{\circ}$} کی قدر کیا ہوگی:}{1}{$\dfrac{\sqrt{3}}{2}$}{$\dfrac{1}{\sqrt{3}}$}{0}
\umcq{246}{ایک قائمۃ الزاویہ مثلث میں اگر ایک زاویہ \lr{30$^{\circ}$} ہو تو حادہ زاویہ کیا ہوگا:}{60$^{\circ}$}{45$^{\circ}$}{90$^{\circ}$}{120$^{\circ}$}
\umcq{247}{ایک قائمۃ الزاویہ مثلث میں اگر ایک سائیڈ \lr{3} سینٹی میٹر اور دوسری سائیڈ \lr{4} سینٹی میٹر ہو تو وتر کی لمبائی کیا ہوگی:}{5 cm}{6 cm}{7 cm}{8 cm}
\umcq{248}{ایک قائمۃ الزاویہ مثلث میں اگر ایک زاویہ \lr{45$^{\circ}$} ہو تو وتر کی لمبائی کیا ہوگی اگر ایک سائیڈ \lr{5} سینٹی میٹر ہو تو:}{$5\sqrt{2}$ cm}{5 cm}{10 cm}{$10\sqrt{2}$}
\umcq{249}{ایک قائمۃ الزاویہ مثلث میں اگر وتر \lr{10} سینٹی میٹر اور ایک سائیڈ \lr{6} سینٹی میٹر ہو تو دوسری طرف کی لمبائی کیا ہوگی:}{\urdu{4 سینٹی میٹر}}{\urdu{6 سینٹی میٹر}}{\urdu{8 سینٹی میٹر}}{\urdu{2 سینٹی میٹر}}
\umcq{250}{ایک قائمۃ الزاویہ مثلث میں اگر ایک سائیڈ \lr{3} سینٹی میٹر اور دوسری سائیڈ \lr{4} سینٹی میٹر ہو تو وتر کی لمبائی کیا ہوگی:}{\urdu{5 سینٹی میٹر}}{\urdu{6 سینٹی میٹر}}{\urdu{7 سینٹی میٹ}}{\urdu{8 سینٹی میٹر}}
\umcq{251}{ریڈین میں \lr{$\tan^{-1}2$} کی قدر ہے::}{$\dfrac{\pi}{2}$}{$\dfrac{3\pi}{2}$}{0.4636 $\pi$}{0.4636}
\umcq{252}{قائمۃ الزاویہ مثلث میں وتر \lr{13} اکائیاں اور ایک زاویہ \lr{$\theta=30^{\circ}$} ہے۔ مخالف ضلع کی لمبائی ہے::}{\urdu{6.5 اکائیاں}}{\urdu{7.5 اکائیاں}}{\urdu{6 اکائیاں}}{\urdu{5 اکائیاں}}
\umcq{253}{عمارت سے \lr{50 m} دور کھڑا شخص عمارت کی چوٹی کو \lr{45$^{\circ}$} کی بلندی کی زاویے پر دیکھتا ہے عمارت کی اونچائی ہے:}{50 m}{25 m}{35 m}{70 m}
\umcq{254}{\lr{$\sec^{2}\theta-\tan^{2}\theta$}:}{$\sin^{2}\theta$}{1}{$\cos^{2}\theta$}{$\cot^{2}\theta$}
\umcq{255}{اگر \lr{$\sin\theta=\dfrac{3}{5}$} اور \lr{$\theta$} ایک حادہ زاویہ ہو تو \lr{$\cos^{2}\theta$=......}}{$\dfrac{7}{25}$}{$\dfrac{24}{25}$}{$\dfrac{16}{25}$}{$\dfrac{4}{25}$}
\umcq{256}{\lr{$\dfrac{5\pi}{24}$ $rad$ = ...... degree:}}{30$^{\circ}$}{37.5$^{\circ}$}{45$^{\circ}$}{52.5$^{\circ}$}
\umcq{257}{\lr{292.5$^{\circ}$= ......rad}:}{$\dfrac{17\pi}{6}$}{$\dfrac{17\pi}{4}$}{1.6 $\pi$}{1.625 $\pi$}
\umcq{258}{درج ذیل میں سے کون سی ایک درست تکونیاتی اکائی ہے؟:}{$\cos\!\left(\dfrac{\pi}{2}-\theta\right)=\sin\theta$}{$\cos\!\left(\dfrac{\pi}{2}-\theta\right)=\cos\theta$}{$\cos\!\left(\dfrac{\pi}{2}-\theta\right)=\sec\theta$}{$\cos\!\left(\dfrac{\pi}{2}-\theta\right)=\cos$}
\umcq{259}{\lr{Sin 60$^{\circ}$= ......}:}{1}{$\dfrac{1}{2}$}{$\sqrt{(3)^{2}}$}{$\dfrac{\sqrt{3}}{2}$}
\umcq{260}{\lr{$\cos^{2}100\,\pi+\sin^{2}100\,\pi$}:}{1}{2}{3}{4}
\umcq{261}{مستوی میں افقی لائن ظاہر کرتی ہے:}{xox$'$}{yoy$'$}{zoz$'$}{None of these}
\umcq{262}{ایک مستوی میں \lr{yoy$'$} کو....... کہا جاتا ہے:}{\urdu{افقی لائن}}{\urdu{عمودی لائن}}{\urdu{مساوی لائن}}{\urdu{ان میں سے کوئی نہیں}}
\umcq{263}{\lr{Ox} پر واقع نقاط کہلاتے ہیں:}{+ ve}{- ve}{+ ve or - ve}{None of these}
\umcq{264}{\lr{(x, y)} میں اگر \lr{x $>$ 0} اور \lr{y $<$ 0} ہو تو نقاط کس ربع میں ہونگے:}{I}{II}{III}{IV}
\umcq{265}{نقطہ \lr{P(3,- 4)} کس ربع میں واقع ہے:}{I}{II}{III}{IV}
\umcq{266}{نقاط \lr{A(2, 3)} اور \lr{B(5, 7)} میں فاصلہ ہے:}{5}{6}{7}{8}
\umcq{267}{دو باہمی ملنے والے لائنوں \lr{(1,2)} اور \lr{(5,8)} کا وسطی نقطہ ہے:}{(3 ,5)}{(4 ,5)}{(3 ,4)}{(4,6)}
\umcq{268}{اگر \lr{(x, 3)} اور \lr{(4,7)} کے درمیان فاصلہ \lr{5} ہو تو \lr{x} ہے:}{1}{2}{3}{4}
\umcq{269}{نقطہ \lr{(0,5)} کس محور پر واقع ہو گا:}{\urdu{x- محور}}{\urdu{y- محور}}{\urdu{دونوں محوروں پر}}{\urdu{کوئی نہیں}}

\umcq{270}{\lr{(5,2)} اور \lr{(9,8)} سے ملنے والے پل کا وسطی نقطہ ہے:}{(5,7)}{(6,7)}{(5,8)}{(6,8)}
\umcq{271}{\lr{(3,2)} اور \lr{(7,4)} میں سے گزرنے والی لائن کی ڈھلوان ہے:}{1}{2}{3}{4}
\umcq{272}{کسی لائن کا ڈھلوان \lr{1} کے ساتھ جھکاؤ ہے:}{30$^{\circ}$}{45$^{\circ}$}{60$^{\circ}$}{90$^{\circ}$}
\umcq{273}{اگر کسی لائن کی ڈھلوان نامعلوم ہے تو لائن ہوگی:}{\urdu{افقی}}{\urdu{عمودی}}{\urdu{\lr{x}-محور کے مساوی}}{\urdu{\lr{y}-محور کے مساوی}}
\umcq{274}{\lr{x}-محور کے مساوی لائن کی ڈھلوان ہوگی:}{0}{1}{Undefined}{-1}
\umcq{275}{کسی لائن کی ڈھلوان \lr{2} ڈھلوانوں والی لائن کے عمودی ہوتی ہے:}{2}{-2}{$\dfrac{1}{2}$}{$-\dfrac{1}{2}$}
\umcq{276}{\lr{3} ڈھلوانوں اور \lr{y}- قاطع \lr{2} والی لائن کی مساوات ہے:}{y = 3x + 2}{y = 2x + 3}{y = 3x - 2}{y = 2x - 3}
\umcq{277}{\lr{4} ڈھلوانوں کے ساتھ \lr{(2,1)} میں سے گزرنے والی لائن کی مساوات ہوگی:}{y = 4x - 2}{y = 4x + 2}{y = 4x - 6}{y = 4x + 6}
\umcq{278}{\lr{(3,2)} اور \lr{(7,4)} میں گزرنے والی لائن کی مساوات ہوگی:}{y = 2x - 1}{y = 2x + 1}{y = 3x - 1}{y = 3x + 1}
\umcq{279}{لائن جس کا \lr{x}- قاطع \lr{3} اور \lr{y}- قاطع \lr{4} ہو تو اس کی مساوات ہوگی:}{4x+3y=12}{3x+4y=12}{4x$-$3y=12}{3x$-$4y=12}
\umcq{280}{لائن جو \lr{y = 2x + 3} کے متوازی ہو اور \lr{(1,1)} سے گزرے اس کی مساوات ہوگی:}{y = 2x - 1}{y = 2x + 1}{y = 2x - 3}{y = 2x + 3}
\umcq{281}{\lr{(4,3)} میں سے گزرنے والی افقی لائن کی مساوات ہے:}{y = 3}{y = 4}{x = 3}{x = 4}
\umcq{282}{\lr{(6,5)} میں سے گزرنے والی عمودی لائن کی مساوات ہے:}{y = 5}{y = 6}{x = 5}{x = 6}
\umcq{283}{لائن \lr{2x + 3y = 6} کی ڈھلوان ہے:}{$\dfrac{2}{3}$}{$-\dfrac{2}{3}$}{$\dfrac{3}{2}$}{$-\dfrac{3}{2}$}
\umcq{284}{اگر دو لائنیں عمودی ہوں تو ان کی ڈھلوانوں کا مجموعہ ہوگا:}{1}{-1}{0}{2}
\umcq{285}{\lr{x = 3} کے عمودی لائن کی ڈھلوان ہوگی:}{0}{1}{2}{-1}
\umcq{286}{مساوات \lr{2x + 3y = 6} ڈھلوان قاطع میں ہوگی:}{$y=\dfrac{-2}{3x+2}$}{$y=\dfrac{2}{3x+2}$}{$y=\dfrac{-2}{3x-2}$}{$y=\dfrac{2}{3x-2}$}
\umcq{287}{مساوات \lr{3x - 4y = 12} قاطع کی شکل میں ہے:}{$\dfrac{x}{4}+\dfrac{y}{3}=1$}{$\dfrac{x}{4}-\dfrac{y}{3}=1$}{$\dfrac{x}{3}+\dfrac{y}{4}=1$}{$\dfrac{x}{3}-\dfrac{y}{4}=1$}
\umcq{288}{مساوات \lr{5x - 12y + 39 = 0} نارمل شکل میں ہے:}{5x $-$ 12y = $-$ 39}{5x + 12y = 39}{5x $-$ 12y = 39}{5x + 12y = $-$ 39}
\umcq{289}{محدودات \lr{(4,3)} اور \lr{(1,7)} والے دو قصبوں کے درمیان فاصلہ ہے:}{5}{6}{7}{8}
\umcq{290}{\lr{A(2, 3) B(5, 7)} اور \lr{C(6,2)} اس والی مثلثی ایک طرف کی لمبائی ہے:}{5}{6}{7}{8}
\umcq{291}{محدودات \lr{(20,10)} اور \lr{(40,30)} والے دو شہروں کا فاصلہ ہوگا:}{$20\sqrt{2}$}{$30\sqrt{2}$}{$40\sqrt{2}$}{$50\sqrt{2}$}
\umcq{292}{نقشے پر محدودات \lr{(5,12)} اور \lr{(8,-4)} والے دو نقاط کا فاصلہ ہے:}{9}{10}{11}{12}
\umcq{293}{\lr{(0,0),(0,5),(8,5)} اور \lr{(8,0)} کونوں والے ایک مستطیل میدان کے گرد جنگلے کی لمبائی ہوگی:}{40}{30}{36}{26}
\umcq{294}{\lr{(40$^{o}$N, 100$^{o}$W)} سے \lr{(50$^{o}$N, 80$^{o}$W)} تک سفر کا وسطی نقطہ ہے:}{(50$^{o}$N, 90$^{o}$W)}{(45$^{o}$N, 80$^{o}$W)}{(45$^{o}$N, 90$^{o}$W)}{(40$^{o}$N, 90$^{o}$W)}
\umcq{295}{ڈھلوان قاطع شکل میں خط مستقیم کی مساوات اس طرح لکھی جاتی ہے:}{y=m(x+c)}{y-y$_1$=m(x$-$x$_1$)}{y = c + mx}{ax + by + c = 0}
\umcq{296}{دو متوازی خطوط کے میلان ہوتے ہیں:}{\urdu{برابر}}{\urdu{صفر}}{\urdu{ایک دوسرے کے منفی رد عمل}}{\urdu{غیر واضح}}
\umcq{297}{اگر دو خطوط کی ڈھلوان کی ضرب \lr{-1} ہو تو خطوط ہوتے ہیں:}{\urdu{متوازی}}{\urdu{عمودی}}{\urdu{ہم خط}}{\urdu{ایک جیسے}}
\umcq{298}{نقاط \lr{P(1,2)} اور \lr{Q(4,6)} کے درمیان فاصلہ ہوتا ہے:}{5}{6}{$\sqrt{13}$}{4}
\umcq{299}{قطعہ خط کے سروں \lr{(-2,4)} اور \lr{(6,-2)} کا وسطی نقطہ ہوتا ہے:}{(4,2)}{(2,1)}{(1,1)}{(0,0)}
\umcq{300}{نقاط \lr{(1,2)} اور \lr{(4,5)} میں سے گزرنے والا خط ہے:}{y = x + 1}{y = 2x + 3}{y = 3x - 2}{y = x + 2}
\umcq{301}{نقطہ ڈھلوان شکل میں خط کی مساوات ہے:}{y=m(x+c)}{y-y$_1$=m(x$-$x$_1$)}{y=c+mx}{ax+by+c=0}

\umcq{302}{\lr{2x+3y$-$6=0} کی ڈھلوان قاطع شکل ہے:}{$y=\dfrac{-2}{3}x+2$}{$y=\dfrac{2}{3}x-2$}{$y=\dfrac{2}{3}x+1$}{$y=\dfrac{-2}{3}x-2$}
\umcq{303}{سمیٹرک شکل میں خط کی مساوات ہے:}{$\dfrac{x}{a}+\dfrac{y}{b}=1$}{$\dfrac{x-x_1}{1}+\dfrac{y-y_1}{m}=\dfrac{z-z_1}{1}$}{ax + by + c = 0}{$y-y_1=m(x-x_1)$}
\umcq{304}{عام شکل میں خط کی مساوات ہے:}{y = mx + c}{$\dfrac{x}{a}+\dfrac{y}{b}=1$}{$\dfrac{x-x_1}{\cos\alpha}=\dfrac{y-y_1}{\sin\alpha}$}{$x\cos\alpha+y\sin\alpha=p$}
\umcq{305}{درج ذیل میں سے کون سا بیان اکثر استقرائی استدلال سے متعلق ہوتا ہے:}{\urdu{بار بار کیے گئے تجربات پر مبنی}}{\urdu{اگر اور صرف اگر کے بیانات}}{\urdu{بیان کو ایک مسئلے کے ذریعے ثابت کیا گیا ہو}}{\urdu{عمومی اصولوں پر مبنی}}
\umcq{306}{درج ذیل میں سے کون سا جملہ استخراجی استدلال کی وضاحت کر تا ہے:}{\urdu{محدود مشاہدات سے عمومی نتائج}}{\urdu{بار بار کیے گئے تجربات پر مبنی}}{\urdu{معلومات کی اکائیوں پر مبنی جو درست ہوں}}{\urdu{معروف حقائق سے نتیجہ اخذ کرنا}}
\umcq{307}{درج ذیل میں سے کون سا بیان درست ہے:}{\urdu{صحیح اعداد کا سیٹ متناہی ہوتا ہے}}{\urdu{کسی بھی چو کور کے اندرونی زاویوں کا مجموعہ ہمیشہ \lr{180$^{o}$} ہوتا ہے}}{\lr{$\dfrac{22}{7}\notin Q'$}}{\urdu{تمام مساوی الاضلاع مثلثیں مساوی الساقین ہوتی ہیں}}
\umcq{308}{چولہا جل رہا ہے کہ بیان کی نفی کے لیے درج ذیل میں سے کون سا بیان بہترین ہے:}{\urdu{چولہا نہیں جل رہا}}{\urdu{چولہا مدھم ہے}}{\urdu{چولہا ہلکی آنچ پر بدل جاتا ہے}}{\urdu{یہ جل رہا ہے اور نہیں جل رہا ہے}}
\umcq{309}{دونوں بیانات \lr{p} اور \lr{q} کنجکشن اس وقت درست ہوتا ہے جب:}{\urdu{\lr{p} اور \lr{q} دونوں غلط ہوں}}{\urdu{\lr{p} اور \lr{q} دونوں درست ہوں}}{\urdu{صرف \lr{q} درست ہو}}{\urdu{صرف \lr{p} درست ہو}}
\umcq{310}{ایک مشروط کو صرف اس صورت میں غلط سمجھا جاتا ہے جب:}{\urdu{مفروضہ درست ہو اور نتیجہ غلط ہو}}{\urdu{نتیجہ درست ہو اور مفروضہ غلط ہو}}{\urdu{صرف مفروضہ درست ہو}}{\urdu{صرف نتیجہ غلط ہو}}
\umcq{311}{\lr{$p\rightarrow q$} کا ضد مثبت برابر ہے::}{$q\rightarrow\sim p$}{$\sim q\rightarrow p$}{$\sim p\rightarrow\sim q$}{$\sim q\rightarrow\sim p$}
\umcq{312}{بیان \lr{2} سے بڑا ہر صحیح عدد دو مفرد اعداد کا مجموعہ ہے یہ ہے:}{\urdu{مسئلہ}}{\urdu{قیاس}}{\urdu{اصول متعارفہ}}{\urdu{منطق}}
\umcq{313}{بیان کسی بھی دو نقاط سے ایک خط کھنچا جا سکتا ہے یہ ہے:}{\urdu{مسئلہ}}{\urdu{قیاس}}{\urdu{اصول متعارفہ}}{\urdu{منطق}}
\umcq{314}{بیان ایک مثلث کے اندرونی زاویوں کا مجموعہ \lr{180$^{o}$} ہوتا ہے یہ ہے:}{\urdu{الٹ}}{\urdu{مسئلہ}}{\urdu{اصول متعارفہ}}{\urdu{مشروط}}
\umcq{315}{ایک جملہ یا ریاضیاتی مساوات جو کبھی غلط ہوتی ہے کبھی درست ہوتی ہے لیکن دونوں نہیں ہو سکتی، کہلاتا ہے:}{\urdu{آرگومنٹ}}{\urdu{بیان}}{\urdu{ثبوت}}{\urdu{پڑتال}}
\umcq{316}{کسی مثلث کے اندرونی زاویوں کا مجموعہ ہے:}{90$^{o}$}{180$^{o}$}{270$^{o}$}{360$^{o}$}
\umcq{317}{\lr{r} ردائس والے دائرے کا محیط ہوگا:}{$2\pi r$}{$\pi r^{2}$}{$4\pi r^{2}$}{None of these}
\umcq{318}{"یا" کے لیے علامت استعمال ہوتی ہے:}{$\wedge$}{$\vee$}{$\sim$}{$\equiv$}
\umcq{319}{علامت \lr{$\wedge$} کس کے لیے استعمال ہوتی ہے::}{}{}{}{}
\umcq{320}{\lr{q} کی نفی ایسے لکھا جائے گا:}{$\sim q$}{$\wedge q$}{$\vee q$}{$\Rightarrow q$}
\umcq{321}{دو بیانوں \lr{p} اور \lr{q} کے درمیان کنجکشن کو علامتی طور پر لکھا جاتا ہے:}{$p\rightarrow\sim q$}{$p\vee q$}{$p\wedge q$}{$p\rightarrow q$}
\umcq{322}{کنجکشن کی مثال ہے:}{\lr{8>7 $\wedge$ 5>3}}{\lr{10>10 $\wedge$ 3>3}}{\lr{6>9 $\wedge$ 5>2}}{None of these}
\umcq{323}{مثلث کے دو قائمۃ الزاویہ ہوں یا لاہور سندھ کا دارلحکومت ہے دیا گیا بیان کہلاتا ہے:}{\urdu{کنجکشن}}{\urdu{ڈس جنکشن}}{\urdu{نفی}}{\urdu{ان میں سے کوئی نہیں}}
\umcq{324}{\lr{$p\rightarrow q$} بنانے کے لیے کمپاؤنڈ بیان کیا کہلاتا ہے:}{\urdu{کنجکشن}}{\urdu{ڈس جنکشن}}{\urdu{نفی}}{\urdu{مشروط}}
\umcq{325}{\lr{$p\rightarrow q$} میں \lr{p} کہلاتا ہے::}{\urdu{مفروضہ}}{\urdu{نتیجہ}}{\urdu{دونوں آن، اور}}{\urdu{ان میں سے کوئی نہیں}}
\umcq{326}{ایک مشروط کو ....... سمجھا جاتا ہے جب سابقہ صحیح ہے اور نتیجہ غلط ہے:}{\urdu{درست}}{\urdu{غلط}}{\urdu{نگیشن}}{\urdu{کوئی نہیں}}

\umcq{327}{بیان \lr{$p\rightarrow q$} کہلاتا ہے:}{\urdu{مشروط}}{\urdu{دو شرطی}}{\urdu{کنجکشن}}{\urdu{ڈس جنکشن}}
\umcq{328}{\lr{$p\rightarrow q$} شرط کی بجائے الٹ ہے:}{$q\rightarrow p$}{$\sim p\rightarrow\sim q$}{$\sim q\rightarrow\sim p$}{None of these}
\umcq{329}{\lr{$p\rightarrow q$}، \lr{$\sim q\rightarrow\sim p$} کی بجائے مشروط ہے، کہلاتی ہے:}{\urdu{انورس}}{\urdu{الٹ}}{\urdu{کانٹ پوزیٹو}}{\urdu{ان میں سے کوئی نہیں}}
\umcq{330}{ایک بیان جو متغیر کی تمام ممکنہ قدروں میں شامل ہے وہ ہے:}{\urdu{ٹاٹولوجی}}{\urdu{ابسورٹیٹی}}{\urdu{کانٹی جینسی}}{\urdu{کوئی نہیں}}
\umcq{331}{ایک بیان جو تمام ممکنہ اقدار کو غلط ہے:}{\urdu{ٹاٹولوجی}}{\urdu{ابسورٹی}}{\urdu{کانٹی جینسی}}{\urdu{ان میں سے کوئی نہیں}}
\umcq{332}{\lr{$\dfrac{a}{b}+\dfrac{c}{d}=\ldots\ldots$}}{$\dfrac{ac}{bd}$}{$\dfrac{a+c}{b+d}$}{$\dfrac{ad+bc}{bd}$}{$\dfrac{ab+cd}{bd}$}
\umcq{333}{\lr{(x$-$1)$^{2}$ + (x+1)$^{2}$ = ...................}:}{4(x$^{2}$+1)}{3(x$^{2}$+1)}{2(x$^{2}$+1)}{(x$^{2}$+1)}
\umcq{334}{زیادہ اطراف سے بند شکل کو کہا جاتا ہے:}{\urdu{دائرہ}}{\urdu{مستطیل}}{\urdu{مثلث}}{\urdu{کثیر الاضلاع}}
\umcq{335}{...... شکلوں کی بناوٹ ایک جیسی مگر سائز مختلف ہوتا ہے:}{\urdu{متماثل}}{\urdu{مماثل}}{\urdu{برابر}}{\urdu{ان میں سے کوئی نہیں}}
\umcq{336}{مماثل کی مثال \lr{......} ہے:}{\urdu{اسکیلین مثلث}}{\urdu{ایکیوٹ زاویہ مثلث}}{\urdu{مساوی مثلث}}{\urdu{آبٹس زاویہ مثلث}}
\umcq{337}{مماثلت کی علامت ہے:}{$\equiv$}{$\sim$}{=}{$\neq$}
\umcq{338}{\lr{ABC} اور \lr{DEF} مماثل ہوں تو ان کو علامتی طور پر لکھا جاتا ہے:}{$\triangle ABC\equiv\triangle DEF$}{$\triangle ABC\sim\triangle DEF$}{$\triangle ABC>\triangle DEF$}{$\triangle ABC<\triangle DEF$}
\umcq{339}{متوازی الاضلاع کا علاقہ:}{\urdu{قاعدہ + اونچائی}}{\urdu{قاعدہ \lr{$\times$} اونچائی}}{\urdu{\lr{$\dfrac{1}{2}$} قاعدہ \lr{$\times$} اونچائی}}{\urdu{لمبائی \lr{$\times$} چوڑائی}}
\umcq{340}{دو کثیر اضلاع کے مماثل ہونے کی شرط کیا ہوگی:}{\urdu{اگر متناظر زوایے اور متناظر اضلاع برابر ہوں}}{\urdu{متناظر زوایے برابر ہوں اور متناظر اضلاع پروپوشنل ہوں}}{\urdu{متناظر زوایے پروپوشنل ہوں اور متناظر اضلاع برابر ہوں}}{\urdu{متناظر زوایے اور متناظر اضلاع پروپوشنل ہوں}}
\umcq{341}{اگر دو مثلثوں کے متناظر اضلاع کی نسبت \lr{3:4} ہو تو ان کے اضلاع کی نسبت کیا ہوگی:}{3:4}{9:8}{9:16}{12:16}
\umcq{342}{اگر مثلث \lr{ABC} اور \lr{DEF} میں اگر \lr{AB =6} سینٹی میٹر اور \lr{DE =9} سینٹی میٹر اور \lr{AC =8} سینٹی میٹر \lr{DF} کیا ہوگا اگر مثلثیں مماثل ہوں:}{\urdu{10 سینٹی میٹر}}{\urdu{12 سینٹی میٹر}}{\urdu{14 سینٹی میٹر}}{\urdu{16 سینٹی میٹر}}
\umcq{343}{اگر قاعدہ \lr{8} سینٹی میٹر اور اونچائی \lr{4} سینٹی میٹر ہو تو متوازی الاضلاع کا رقبہ ہوگا:}{12 cm$^{2}$}{4 cm$^{2}$}{18 cm$^{2}$}{32 cm$^{2}$}
\umcq{344}{اگر ایک چھوٹی کثیر اضلاع کا رقبہ \lr{25cm$^{2}$} ہو اور متناظر اضلاع کی نسبت \lr{1:2} ہو تو، بڑی کثیر اضلاع کا رقبہ کیا ہوگا:}{100 cm$^{2}$}{75 cm$^{2}$}{50 cm$^{2}$}{125 cm$^{2}$}
\umcq{345}{کثیر اضلاع کے بیرونی زاویوں کا فارمولا ہے:}{$\dfrac{360^{o}}{n}$}{$\dfrac{180^{o}}{n}$}{$\dfrac{(n-2)\times180^{o}}{n}$}{None of these}
\umcq{346}{با قاعدہ کثیر الاضلاع کے اندرونی زاویوں کا سائز ہوتا ہے:}{$\dfrac{(n-2)\times180^{o}}{n}$}{$\dfrac{180^{o}}{n}$}{$\dfrac{360^{o}}{n}$}{None of these}
\umcq{347}{مسدس کے اندرونی زاویوں کا مجموعہ کیا ہوگا:}{720$^{o}$}{540$^{o}$}{900$^{o}$}{1080$^{o}$}
\umcq{348}{منظم مخمس کے ہر اندرونی زوایے کی پیمائش کیا ہوگی:}{72$^{o}$}{108$^{o}$}{120$^{o}$}{144$^{o}$}
\umcq{349}{منظم مثمن کے ہر بیرونی زوایے کی پیمائش کیا ہوگی:}{30$^{o}$}{90$^{o}$}{60$^{o}$}{45$^{o}$}
\umcq{350}{ایک منظم سات اضلاع والی شکل کے کتنے وتر ہوتے ہیں:}{7}{21}{14}{28}
\umcq{351}{کسی بھی کثیر اضلاع کے بیرونی زاویوں کا مجموعہ کیا ہوتا ہے:}{180$^{o}$}{360$^{o}$}{540$^{o}$}{720$^{o}$}
\umcq{352}{ایک مساوی الاضلاع مثلث میں، ہر زاویے کی پیمائش کیا ہوگی:}{45$^{o}$}{60$^{o}$}{90$^{o}$}{120$^{o}$}
\umcq{353}{ایک متساوی اضلاع مثلث اگر ایک زاویہ \lr{80$^{o}$} ہو تو اس کو دو زاویوں کا مجموعہ کیا ہوگا:}{50$^{o}$}{60$^{o}$}{70$^{o}$}{80$^{o}$}

\umcq{354}{ایک متوازی اضلاع میں، دو مخالف اطراف میں کیا تعلق ہوتا ہے:}{\urdu{برابر اور پیرالل}}{\urdu{برابر لیکن پیرالل نہیں ہوتیں}}{\urdu{پیرالل لیکن برابر نہیں ہوتیں}}{\urdu{نہ ہی پیرالل اور نہ ہی برابر ہوتیں ہیں}}
\umcq{355}{ایک مثلث میں ہر زوایے کی پیمائش کیا ہوتی ہے:}{60$^{o}$}{120$^{o}$}{90$^{o}$}{180$^{o}$}
\umcq{356}{ایک رومبس میں، وتروں کے درمیان کیا تعلق ہوتا ہے:}{\urdu{مساوی}}{\urdu{غیر مساوی}}{\urdu{پیرالل}}{\urdu{ان میں سے ہر ایک کو صحیح زاویوں پر قطع کرتے ہیں}}
\umcq{357}{اگر دو کثیر الاضلاع متشابہ ہوں تو:}{\urdu{ان کے متناظر زاویے برابر ہوتے ہیں}}{\urdu{ان کے رقبے برابر ہوتے ہیں}}{\urdu{ان کے حجم برابر ہوتے ہیں}}{\urdu{ان کے متناظر اضلاع برابر ہوتے ہیں}}
\umcq{358}{دو متشابہ کثیر الاضلاع کے رقبوں کی نسبت:}{\urdu{ان کے احاطوں کی نسبت کے برابر ہوتی ہے}}{\urdu{ان کے متناظر اضلاع کی لمبائیوں کے مربع کی نسبت کے برابر ہوتی ہے}}{\urdu{ان کے متناظر اضلاع کی لمبائیوں کے مکعب کی نسبت کے برابر ہوتی ہے}}{\urdu{ان کے متناظر اضلاع کے مجموعے کے برابر ہوتی ہے}}
\umcq{359}{اگر دو متشابہ اجسام کے حجم \lr{125 cm$^{3}$} اور \lr{27 cm$^{3}$} ہوں تو ان کی متناظر بلندی کی نسبت ہے:}{3:5}{5:3}{25:9}{9:25}
\umcq{360}{منظم مخمس کا بیرونی زاویہ ہوتا ہے:}{40$^{o}$}{45$^{o}$}{60$^{o}$}{72$^{o}$}
\umcq{361}{ایک متوازی الاضلاع کا رقبہ \lr{64 cm$^{2}$} ہے اور اس کے متشابہ متوازی اضلاع کا رقبہ \lr{144 cm$^{2}$} ہے۔ اگر چھوٹی متوازی الاضلاع کی لمبائی \lr{8} سینٹی میٹر ہو تو بڑی متوازی الاضلاع کے متناظر ضلع کی لمبائی ہے:}{10 cm}{12 cm}{18 cm}{16 cm}
\umcq{362}{ایک نو اضلاع والی کثیر اضلاع میں کل وتروں کی تعداد ہے:}{18}{21}{25}{27}
\umcq{363}{دو متشابہ کروں کے رداس میں نسبت \lr{4:5} ہے اگر بڑے کرے کی سطح کا رقبہ \lr{500$\pi$ cm$^{2}$} ہو تو چھوٹے کرے کی سطح کا رقبہ ہے::}{256$\pi$ cm$^{2}$}{320$\pi$ cm$^{2}$}{400$\pi$ cm$^{2}$}{405$\pi$ cm$^{2}$}
\umcq{364}{ایک منظم کثیر الاضلاع کا بیرونی زاویہ \lr{30$^{o}$} ہے کثیر اضلاع میں وتروں کی تعداد ہے:}{54}{90}{72}{108}
\umcq{365}{ایک منظم مسدس میں وتر کی لمبائی اور ضلع کی لمبائی میں نسبت ہے:}{$\sqrt{3}:1$}{2:1}{3:2}{2:3}
\umcq{366}{ایک منظم کثیر الاضلاع کا اندرونی زاویہ \lr{165$^{o}$} ہے اس کے اضلاع کی تعداد ہے:}{15}{16}{20}{24}
\umcq{367}{یک درجی تفاعل کی عام شکل کیا ہے:}{$f(x)=mx+c$}{$f(x)=ax^{2}+bx+c$}{$f(x)=ax^{3}+bx^{2}+cx+d$}{$f(x)=k$}
\umcq{368}{یک درجی تفاعل \lr{y=3x$-$5} میں ڈھلوان \lr{m} کیا ہے:}{5}{3}{-3}{-5}
\umcq{369}{یک درجی تفاعل \lr{y= $-$2x+8} کا \lr{y}- قاطع کیا ہے:}{-2}{2}{-8}{8}
\umcq{370}{\lr{y=x$^{2}$} کا گراف ظاہر کرتا ہے:}{\urdu{ایک سیدھی لائن}}{\urdu{اوپر کی جانب کھلنے والا پیرابولا}}{\urdu{نیچے کی طرف کھلنے والا پیرابولا}}{\urdu{ایک کیوبک کرو}}
\umcq{371}{دو درجی تفاعل \lr{y=$-$x$^{3}$} میں پیرابولا کی سمت کیا ہوگی:}{\urdu{اوپر کی طرف}}{\urdu{نیچے کی طرف}}{\urdu{بائیں}}{\urdu{دائیں}}
\umcq{372}{سہ درجی تفاعل کی معیاری شکل ہوتی ہے:}{y=ax$^{3}$+bx$^{2}$+cx+d}{y=ax$^{2}$+bx+c}{y=ax+b}{None of these}
\umcq{373}{\lr{y=2x$^{3}$-3x$^{2}$+x-5} کا گراف کس نقطے سے گزرتا ہے:}{(0,-5)}{(0,0)}{(5,0)}{(0,-10)}
\umcq{374}{معکوس تفاعل کی عام شکل ہوتی ہے:}{y=ax+b}{y=ax$^{2}$+bx+c}{$y=\dfrac{q}{x}$}{y=x$^{3}$}
\umcq{375}{\lr{(}ایکویشن\lr{)} کس قسم کا تفاعل ہے:}{\urdu{قوت نمائی}}{\urdu{یک درجی}}{\urdu{دو درجی}}{\urdu{سہ درجی}}
\umcq{376}{\lr{y=x$^{3}$} کا گراف کیسا دیکھائی دیتا ہے:}{\urdu{ایک سیدھی لائن}}{\urdu{ایک پیرابولا}}{\urdu{ایک کیوبک کروس مرکز سے گزرتا ہوا}}{\urdu{ایک ہاپربولا}}
\umcq{377}{تفاعل \lr{y=5x} کا گراف ظاہر کرتا ہے:}{\urdu{قوت نمائی بڑھوتری}}{\urdu{قوت نمائی انخطاط}}{\urdu{ایک لائن مرکز سے گزر رہی ہے}}{\urdu{ایک ہائپربولا}}
\umcq{378}{دو درجی تفاعل \lr{y = x$^{2}$+ 2x - 3} کے لیے بنیادی نقطہ کیا ہوگا:}{(-1,-3)}{(0 ,3)}{(1 ,2)}{(-2 ,1)}
\umcq{379}{\lr{y = - x$^{2}$+ 4} کا گراف \lr{x} قاطع کو قطع کرتا ہے:}{(0 ,4)}{(0,-4)}{(4 ,0)}{(-4 ,0)}

\umcq{380}{لائن \lr{y = 2x + 3} کی ڈھلوان کیا ہوگی:}{2}{3}{-2}{0}
\umcq{381}{لائن \lr{y = -3x + 2} کی ڈھلوان کیا ہوگی:}{3}{-3}{2}{-2}
\umcq{382}{مساوات \lr{y = 3x} کس قسم کے گراف کو ظاہر کرتی ہے:}{\urdu{ایک \lr{3} ڈھلوانوں والی لائن}}{\urdu{ایک پیرابولا}}{\urdu{ایک منحنی خط}}{\urdu{قوت نمائی انخطاط}}
\umcq{383}{\lr{x = 1} پر منحنی \lr{y = 2x$^{2}$- 3x + 1} کی طرف مماس کی ڈھلوان کیا ہوگی:}{1}{2}{3}{4}
\umcq{384}{مساوات \lr{y = - 2x$^{2}$+ 4x - 3} کس قسم کے گراف کو ظاہر کرتی ہے:}{\urdu{پیرابولا نیچے کی طرف کھلتا ہے}}{\urdu{پیرابولا اوپر کی طرف کھلتا ہے}}{\urdu{لائن}}{\urdu{قوت نمائی بڑھوتری}}
\umcq{385}{ایک معکوس تناسب \lr{$y=\dfrac{1}{x}$} میں جیسے \lr{x} صفر ہو جاتا ہے تو \lr{y} کے ساتھ کیا ہوتا ہے:}{\urdu{یہ لامتناہی تک پہنچ جاتا ہے}}{\urdu{یہ صفر پر پہنچ جاتا ہے}}{\urdu{یہ ناقابل بیان ہو جاتا ہے}}{\urdu{یہ پہلے جیسا رہتا ہے}}
\umcq{386}{دو درجی تفاعل کی عام شکل ہے:}{y = mx + c}{y = ax$^{2}$+ bx + c}{y = ax$^{3}$+ bx$^{2}$+ cx + d}{y = k x e$^{2}$x}
\umcq{387}{تفاعل \lr{y = x$^{2}$ + 4x + 4} میں \lr{y} کی قدر کیا ہوگی جب \lr{x = -2} ہو تو:}{0}{4}{-4}{8}
\umcq{388}{مساوات \lr{y = 5x - 2} میں \lr{y} قطع کیا ہے:}{-5}{-2}{2}{5}
\umcq{389}{\lr{C(x) = 1200 + 20x} اور \lr{R(x) = 50x} ہو تو مساوی نقطہ کیا ہوگا:}{40}{50}{60}{75}
\umcq{390}{تفاعل \lr{y = x$^{3}$ + 3x$^{2}$} میں \lr{x = 0} پر مماس کی ڈھلوان کیا ہوگی:}{0}{3}{6}{-3}
\umcq{391}{تفاعل \lr{y = 3x$^{2}$+ 2x - 1} اپنی کم سے کم قدر کس نقطہ پر رکھتا ہے:}{x = 0}{x = $-$1}{x = 1}{x = 2}
\umcq{392}{تفاعل \lr{y = x$^{2}$ + 4x} میں سمیٹری کا قاطع کیا ہوگا:}{x = 0}{x = $-$2}{x = 2}{x = $-$4}
\umcq{393}{درج ذیل میں سے معکوس تفاعل (ایکویشن) کی ڈومین کیا ہے:}{\urdu{تمام حقیقی اعداد سوائے صفر کے}}{\urdu{تمام حقیقی اعداد}}{\urdu{مثبت حقیقی اعداد}}{\urdu{منفی حقیقی اعداد}}
\umcq{394}{\lr{y = 2x$^{3}$- x$^{2}$ - 5} کا گراف ہے:}{\urdu{یک درجی}}{\urdu{دو درجی}}{\urdu{سہ درجی}}{\urdu{قوت نمائی}}
\umcq{395}{\lr{y = 2x$^{3}$- x$^{2}$ - 5} کا گراف ہے:}{\urdu{یک درجی}}{\urdu{دو درجی}}{\urdu{سہ درجی}}{\urdu{قوت نمائی}}
\umcq{396}{تفاعل \lr{y = 3x - 4} کے لیے گراف ہے:}{\urdu{لائن}}{\urdu{پیرابولا}}{\urdu{ہائپربولا}}{\urdu{قوت نمائی تفاعل}}
\umcq{397}{\lr{y = x$^{3}$} کا گراف کس نقطے سے گزرتا ہے:}{(0,0)}{(1,0)}{(0,1)}{(0,-1)}
\umcq{398}{مساوات \lr{y = 4x - 6} کا \lr{y} قاطع کیا ہوگا:}{-6}{6}{4}{-4}
\umcq{399}{کوسٹ تفاعل \lr{C(x) = 1500 + 10x} اور ریوینیو تفاعل \lr{R(x) = 50x} ہو تو مساوی نقطہ کیا ہوگا:}{\urdu{30 یونٹس}}{\urdu{50 یونٹس}}{\urdu{100 یونٹس}}{\urdu{150 یونٹس}}
\umcq{400}{گراف \lr{y = x$^{2}$ - 2x + 1} ظاہر کرتا ہے:}{\urdu{لائن}}{\urdu{پیرابولا}}{\urdu{منحنی خط}}{\urdu{قوت نمائی خط}}
\umcq{401}{تفاعل \lr{y = 2x$^{2}$ + 3x + 5} میں \lr{y} کی قدر کیا ہوگی، جب \lr{x = -1}:}{4}{2}{1}{3}
\umcq{402}{\lr{y = x$^{2}$ - 4} کا گراف \lr{x} قاطع کو کس نقطے پر قطع کرتا ہے:}{\lr{(0,2)} \urdu{اور} \lr{(0,-2)}}{\lr{(4,0)} \urdu{اور} \lr{(-4,0)}}{\lr{(1,0)} \urdu{اور} \lr{(-1,0)}}{\lr{(0,3)} \urdu{اور} \lr{(0,-3)}}
\umcq{403}{تفاعل \lr{y = x$^{4}$} ہے:}{\urdu{یک درجی تفاعل}}{\urdu{دو درجی تفاعل}}{\urdu{سہ درجی تفاعل}}{\urdu{کو آرٹیک تفاعل}}
\umcq{404}{\lr{y = x$^{2}$ + x - 2} کا گراف ظاہر کرتا ہے:}{\urdu{ایک سیدھی لائن}}{\urdu{ہائپر بولا}}{\urdu{پیرابولا}}{\urdu{قوت نمائی تفاعل}}
\umcq{405}{تفاعل \lr{y = x$^{2}$ + 2x + 1} کے لیے پیرابولا کا بنیادی نقطہ کیا ہوگا:}{(0,0)}{(-1,0)}{(-1,1)}{(1,1)}
\umcq{406}{\lr{y = x$^{2}$ + 5} کا گراف \lr{x} قاطع کو کس نقطے پر قطع کرتا ہے:}{(0,5)}{(5,0)}{(-5,0)}{(1,0)}
\umcq{407}{تفاعل \lr{y = 3x$^{2}$} ظاہر کرتا ہے:}{\urdu{لائن کو}}{\urdu{پیرابولا اوپر کی جانب کھلتا ہے}}{\urdu{پیرابولا نیچے کی جانب کھلتا ہے}}{\urdu{سہ درجی تفاعل}}
\umcq{408}{تفاعل \lr{y = 3x$^{2}$} ظاہر کرتا ہے:}{\urdu{ایک لائن}}{\urdu{پیرابولا اوپر کی جانب کھلتا ہے}}{\urdu{پیرابولا نیچے کی جانب کھلتا ہے}}{\urdu{ایک کیوبک تفاعل}}
\umcq{409}{تفاعل \lr{y = 2$^{x}$} کس قسم کی منحنی کو ظاہر کرتا ہے:}{\urdu{یک درجی}}{\urdu{قوت نمائی}}{\urdu{قوت نمائی گراوٹ}}{\urdu{دو درجی}}
\umcq{410}{\lr{x=1} پر \lr{y=x$^{2}$} منحنی خط کی ڈھلوان کیا ہوگی:}{4}{2}{0}{8}

\umcq{411}{کون سا تفاعل قوت نمائی انخطاط کو ظاہر کرتا ہے:}{y=2$^{x}$}{y=e$^{x}$}{y=2$^{-x}$}{y=x$^{2}$}
\umcq{412}{تفاعل \lr{y=x$^{2}$-3x-2} میں \lr{y} کی قدر کیا ہوگی جبکہ \lr{x=0}:}{-2}{0}{2}{-3}
\umcq{413}{تفاعل \lr{y=x$^{2}$-2} کا گراف ظاہر کرتا ہے:}{\urdu{پیرابولا اوپر کی جانب کھلتا ہے}}{\urdu{پیرابولا نیچے کی جانب کھلتا ہے}}{\urdu{ہائپربولا}}{\urdu{دائرہ}}
\umcq{414}{جب کوسٹ تفاعل \lr{C(x)=90,000+600x} اور ریوینیو تفاعل \lr{R(x)=1,800x} ہو تو مساوی نقطہ کیا ہوگا:}{\urdu{50 فٹ بال}}{\urdu{75 فٹ بال}}{\urdu{100 فٹ بال}}{\urdu{150 فٹ بال}}
\umcq{415}{درج ذیل میں سے کون سا قوت نمائی بڑھوتری کا تفاعل ہے:}{y = x$^{2}$}{y = 3x}{y = 2$^{x}$}{y = x$^{3}$}
\umcq{416}{درج ذیل میں سے کون سا قوت نمائی انخطاط تفاعل کو ظاہر کرتا ہے:}{y = 2$^{x}$}{y = e$^{-x}$}{y = x$^{2}$}{y = 3x}
\umcq{417}{\lr{y = e$^{x}$} کا گراف ظاہر کرتا ہے:}{\urdu{یک درجی بڑھوتری}}{\urdu{قوت نمائی بڑھوتری}}{\urdu{پیرابولک بڑھوتری}}{\urdu{سہ درجی بڑھوتری}}
\umcq{418}{\lr{x = 5} ظاہر کرتا ہے:}{\urdu{\lr{x}- محور}}{\urdu{\lr{y}- محور}}{\urdu{\lr{x}- محور کے متوازی خط}}{\urdu{\lr{y}- محور کے متوازی خط}}
\umcq{419}{خط \lr{y = 5x + 3} کی ڈھلوان ہے:}{3}{-3}{5}{-5}
\umcq{420}{\lr{y= -2x - 1} کا \lr{y} قاطع ہے:}{-2}{2}{-1}{1}
\umcq{421}{\lr{$y=x^{3}$} کا گراف \lr{x}- محور کو ...... پر کاٹتا ہے::}{x=0}{x=1}{x=-1}{x=2}
\umcq{422}{\lr{$3^{x}$} کا گراف ظاہر کرتا ہے:}{\urdu{افزائش}}{\urdu{انخطاط}}{\urdu{الف اور ب دونوں}}{\urdu{خط کو}}
\umcq{423}{\lr{$y=-x^{2}+5$} کا گراف کھلتا ہے::}{\urdu{اوپر کی طرف}}{\urdu{نیچے کی طرف}}{\urdu{بائیں طرف}}{\urdu{دائیں طرف}}
\umcq{424}{\lr{$y=x^{2}-9$} کا گراف کھلتا ہے::}{\urdu{اوپر کی طرف}}{\urdu{نیچے کی طرف}}{\urdu{بائیں طرف}}{\urdu{دائیں طرف}}
\umcq{425}{\lr{y=5$^{x}$} ...... تفاعل ہے:}{\urdu{یک درجی}}{\urdu{دو درجی}}{\urdu{سہ درجی}}{\urdu{قوت نمائی}}
\umcq{426}{معکوس تفاعل ہے:}{$y=7^{x}$}{$y=\dfrac{2}{x}$}{$y=2x^{2}$}{$y=5x^{3}$}
\umcq{427}{\lr{$y=-3x^{3}+7$} ...... تفاعل ہے:}{\urdu{قوت نمائی}}{\urdu{سہ درجی}}{\urdu{یک درجی}}{\urdu{معکوس}}
\umcq{428}{مثلث کے اطراف کے عمود ناصفوں کے قطع کا نقطہ کہلاتا ہے:}{\urdu{محور مرکز}}{\urdu{مرکز نما}}{\urdu{محاصرہ مرکز}}{\urdu{عمودی مرکز}}
\umcq{429}{اگر کسی مثلث کی پیمائش کا ایک زاویہ \lr{90$^{\circ}$} سے زیادہ ہو تو کون سی مثلث بنے گی:}{\urdu{حادہ مثلث}}{\urdu{قائمۃ الزاویہ مثلث}}{\urdu{منفرجہ مثلث}}{\urdu{متساوی الساقین مثلث}}
\umcq{430}{وہ نقطہ جہاں کسی مثلث کے تمام وسطانیے ملتے ہیں کہلاتی ہے:}{\urdu{محاصرہ مرکز}}{\urdu{محصور مرکز}}{\urdu{مرکز نما}}{\urdu{عمودی مرکز}}
\umcq{431}{کون سا لوکس کس نقطہ کی وضاحت کرتا ہے جو ایک مقررہ نقطہ سے ایک مقررہ فاصلہ ہے:}{\urdu{عمودی ناصف}}{\urdu{دائرہ}}{\urdu{زاویہ کا ناصف}}{\urdu{لائن سیگمنٹ}}
\umcq{432}{کسی مثلث کے زاویے کا ناصف کیا ہے:}{\urdu{راس سے گزرتی ہوئی لائن}}{\urdu{لائن جو دو زاویوں کو دو برابر حصوں میں تقسیم کرے}}{\urdu{لائن جو دو اطرافوں کے مرکزی نقطہ کو ملائے}}{\urdu{ایک طرف لگی ہوئی لائن}}
\umcq{433}{کسی مثلث کا عمودی مرکز کہاں ہوتا ہے:}{\urdu{مثلث کے ارتفاع ملتے ہیں}}{\urdu{اطراف عمودی ہوتی ہیں}}{\urdu{محاصرہ مرکز مرکز نما سے ملتا ہے}}{\urdu{ان میں سے کوئی نہیں}}
\umcq{434}{مثلث کے تین زاویہ کے دو حصوں کے قطع کو کہتے ہیں:}{\urdu{عمودی مرکز}}{\urdu{مرکز نما}}{\urdu{محاصرہ مرکز}}{\urdu{محصور مرکز}}
\umcq{435}{ایک مساوی الاضلاع مثلث ...... زوایے والے مثلث ہوتی ہے:}{\urdu{منفرجہ}}{\urdu{قائمہ}}{\urdu{حادہ}}{\urdu{ان میں سے کوئی نہیں}}
\umcq{436}{ایک دائرے کا مرکز جو ایک مثلث کے تمام عمودوں میں سے گزرتا ہے کہاں واقع ہوتا ہے:}{\urdu{محصور مرکز}}{\urdu{محاصر مرکز}}{\urdu{مرکز نما}}{\urdu{عمودی مرکز}}
\umcq{437}{زاویے کے لحاظ سے مثلث کی کتنی اقسام ہیں:}{2}{3}{4}{5}
\umcq{438}{ایسی مثلث جس کے تمام زاویے اور اطراف برابر ہوں کہلاتی ہے:}{\urdu{قائمۃ الزاویہ}}{\urdu{متساوی الساقین مثلث}}{\urdu{مساوی الاضلاع مثلث}}{\urdu{مختلف الاضلاع مثلث}}
\umcq{439}{کسی مثلث کا مرکز نما ...... کے ہم نقطہ کو:}{\urdu{زاویوں کے ناصف}}{\urdu{ارتفاع}}{\urdu{عمودی ناصف}}{\urdu{وسطانیہ}}
\umcq{440}{قائمۃ الزاویہ مثلث بناتے ہوئے، قائمہ الزاویہ کے مخالف زاویے کو کہتے ہیں:}{\urdu{ہائپوٹینیس}}{\urdu{وسطانیہ}}{\urdu{ارتفاع}}{\urdu{قاعدہ}}

\umcq{441}{مساوی الاضلاع مثلث کے حوالے سے کیا درست ہے:}{\urdu{یہ \lr{90$^{\circ}$} کا ایک زاویہ رکھتی ہے}}{\urdu{تمام زاویے \lr{60$^{\circ}$} کے ہیں}}{\urdu{یہ دو مساوی اطراف رکھتی ہے}}{\urdu{تمام اطراف مختلف ہوتی ہیں}}
\umcq{442}{مثلث کے تمام زاویوں کا مجموعہ کیا ہوتا ہے:}{180$^{\circ}$}{360$^{\circ}$}{90$^{\circ}$}{120$^{\circ}$}
\umcq{443}{دو برابر اضلاع والی مثلث کیا کہلاتی ہے:}{\urdu{مساوی الاضلاع مثلث}}{\urdu{مختلف الاضلاع مثلث}}{\urdu{متساوی الساقین مثلث}}{\urdu{قائمۃ الزاویہ مثلث}}
\umcq{444}{اضلاع کے لحاظ سے مثلث کی کتنی اقسام ہیں:}{2}{3}{4}{5}
\umcq{445}{کسی مثلث کا مرکز نما ہر وسطانیہ کو دو حصوں میں کس نسبت میں تقسیم کرتا ہے:}{1:2}{2:1}{1:3}{3:1}
\umcq{446}{اضلاع کے لحاظ سے درج ذیل میں سے کون مثلث کی قسم نہیں ہے:}{\urdu{مختلف الاضلاع}}{\urdu{متساوی الساقین}}{\urdu{مساوی الاضلاع}}{\urdu{حادہ}}
\umcq{447}{مثلث کے راس سے اس کے مخالف ضلع پر عمود کیا کہلاتی ہے:}{\urdu{ارتفاع}}{\urdu{وسطانیہ}}{\urdu{ناصف}}{\urdu{زاویہ}}
\umcq{448}{ایک مثلث جس کا ایک زاویہ \lr{90} ڈگری کے برابر ہو کہلاتی ہے:}{\urdu{مختلف الاضلاع مثلث}}{\urdu{مساوی الاضلاع مثلث}}{\urdu{قائمۃ الزاویہ مثلث}}{\urdu{متساوی الساقین مثلث}}
\umcq{449}{ایک بند مثلث کی ہوتی ہیں:}{\urdu{دو اطراف}}{\urdu{تین اطراف}}{\urdu{چار اطراف}}{\urdu{پانچ اطراف}}
\umcq{450}{ایک مثلث بناتے ہوئے، ایک زاویے اور دو اطراف کی پیمائش کے لیے کون سی شرط درکار ہوتی ہے:}{\urdu{ایس اے ایس}}{\urdu{اے ایس اے}}{\urdu{ایس ایس ایس}}{\urdu{اے اے ایس}}
\umcq{451}{ایک مثلث میں راسوں میں سے ایک دوسرے کو قطع کرنے والی شاعیں کون سا زاویہ بناتی ہیں:}{\urdu{تکمیلی زاویہ}}{\urdu{مخالف زاویہ}}{\urdu{عمودی زاویہ}}{\urdu{بیرونی زاویہ}}
\umcq{452}{ایک مثلث بنائی جا سکتی ہے اگر کسی دو اضلاع کا مجموعہ تیسرے ضلع کی پیمائش سے ...... ہو:}{\urdu{کم}}{\urdu{زیادہ}}{\urdu{برابر}}{\urdu{زیادہ اور برابر}}
\umcq{453}{ایک متساوی اضلاع مثلث ......:}{\urdu{متساوی الساقین ہو سکتی ہے}}{\urdu{قائمۃ الزاویہ ہو سکتی ہے}}{\urdu{منفرجہ زاویہ ہو سکتی ہے}}{\urdu{اس کا ہر زاویہ \lr{50$^{o}$} کا ہوتا ہے}}
\umcq{454}{اگر دونوں زاویوں کا مجموعہ \lr{90$^{o}$} سے کم ہو، تو مثلث ...... ہوتی ہے:}{\urdu{متساوی الاضلاع}}{\urdu{حادہ زاویہ}}{\urdu{منفرجہ زاویہ}}{\urdu{قائمۃ الزاویہ}}
\umcq{455}{وہ قطعہ خط جو ایک مثلث کے ضلع کے وسطی نقطہ کو مخالف راس سے ملاتا ہے، کہلاتا ہے:}{\urdu{وسطانیہ}}{\urdu{عمودی ناصف}}{\urdu{زاویہ کا ناصف}}{\urdu{دائرہ}}
\umcq{456}{مثلث کے ناصف زاویے ...... پر قطع کرتے ہیں:}{\urdu{ایک نقطہ}}{\urdu{دو نقاط}}{\urdu{تین نقاط}}{\urdu{چار نقاط}}
\umcq{457}{تمام نقاط کا لوکس جو کسی معین نقطہ سے برابر فاصلے پر ہو ...... کہلاتا ہے:}{\urdu{دائرہ}}{\urdu{عمودی ناصف}}{\urdu{زاویہ کا ناصف}}{\urdu{متوازی خطوط}}
\umcq{458}{نقاط کا لوسائی جو دو معین نقاط سے برابر فاصلے پر ہوں ...... کہلاتا ہے:}{\urdu{دائرہ}}{\urdu{عمودی ناصف}}{\urdu{زاویہ کا ناصف}}{\urdu{متوازی خطوط}}
\umcq{459}{نقاط کو لوسائی جو کسی معین خط سے برابر فاصلے پر ہوں ...... کہلاتا ہے:}{\urdu{دائرہ}}{\urdu{عمودی ناصف}}{\urdu{زاویہ کا ناصف}}{\urdu{متوازی خطوط}}
\umcq{460}{نقاط کا لوسائی جو دو باہم ملتے خطوط سے برابر فاصلے پر ہوں ...... کہلاتا ہے:}{\urdu{دائرہ}}{\urdu{عمودی ناصف}}{\urdu{زاویہ کا ناصف}}{\urdu{متوازی خطوط}}
\umcq{461}{تمام نقاط کا سیٹ جو کسی معین نقطہ \lr{B} سے \lr{2} کلو میٹر سے زیادہ کے فاصلے پر ہوں ایک دائرے کے باہر کا علاقہ جس کا رداس ...... اور مرکز \lr{B} ہوتا ہے:}{1 km}{1.9 km}{2 km}{2.1 km}
\umcq{462}{اعدادی مواد کو جمع کرنے، ترتیب دینے، تلخیص کرنے، تجزیہ کرنے اور تشریح کرنے کو کہتے ہیں:}{\urdu{مواد اکٹھا کرنا}}{\urdu{معلوماتی معاملات}}{\urdu{مواد کا تجزیہ}}{\urdu{مواد کی تشریح}}
\umcq{463}{درج ذیل میں سے کونسی غیر مسلسل مواد کی مثال ہے:}{\urdu{طلبہ کی اونچائی}}{\urdu{طلبہ کا وزن}}{\urdu{فروخت کی گئی کتابیں}}{\urdu{ایک شہر کا درجہ حرارت}}
\umcq{464}{اعداد و شمار کی اصطلاح کیا ہے جو دیئے گئے وقفے میں ہر ممکنہ قدر لے سکتا ہے:}{\urdu{غیر مسلسل مواد}}{\urdu{مسلسل مواد}}{\urdu{گروہی مواد}}{\urdu{غیر گروہی مواد}}
\umcq{465}{درج ذیل میں سے کون غیر گروہی مواد کی مثال ہے:}{\urdu{کلاسوں میں ترتیب دیئے گئے طلبہ کی نمبر}}{\urdu{ایک مہینے میں فروخت کیے گئے کھلونوں کی تعداد}}{\urdu{ایک کلاس میں طلبہ کی اونچائی}}{\urdu{روزانہ محفوظ کیا گیا درجہ حرارت}}
\umcq{466}{منظم ترتیب (گروپ یا کلاسز) میں ترتیب دیئے گئے ڈیٹا کی اصطلاح کیا ہے:}{\urdu{غیر گروہی مواد}}{\urdu{گروہی مواد}}{\urdu{غیر مسلسل مواد}}{\urdu{مسلسل مواد}}
\umcq{467}{معلوماتی معاملات کا استعمال دنیا میں سب سے پہلے کس نے کیا:}{\urdu{نیوٹن}}{\urdu{ہنری}}{\urdu{سرروندفنشر}}{\urdu{آرتھر کیلی}}
\umcq{468}{جماعتی حدود کو معلوم کرنے کا فارمولا کیا ہے:}{\urdu{سائز = رینج / جماعتوں کی تعداد}}{\urdu{سائز = رینج \lr{x} جماعتوں کی تعداد}}{\urdu{سائز = رینج + جماعتوں کی تعداد}}{\urdu{سائز = رینج - جماعتوں کی تعداد}}
\umcq{469}{اگر سب سے بڑی قدر \lr{136} اور سب سے چھوٹی قدر \lr{30} ہو تو مواد کی سعت کیا ہوگی:}{106}{166}{96}{136}
\umcq{470}{بالائی اور زیریں جماعتی فرق کے لیے کون سی اصلاح استعمال ہوگی:}{\urdu{جماعت کا سائز}}{\urdu{جماعتی حدود}}{\urdu{جماعتی وقفہ}}{\urdu{جماعتی تعداد}}
\umcq{471}{جماعت \lr{10-14} میں زیریں حد کیا ہوگی:}{9.5}{10}{12}{146}

\umcq{472}{کالمی نقشہ کیا ہے:}{\urdu{ملحقہ مستطیلوں کا گراف}}{\urdu{جڑی ہوئی لائنوں کا گراف}}{\urdu{پائی پارٹ کا گراف}}{\urdu{بار چارٹ کا گراف}}
\umcq{473}{جماعتی حدود کے درمیانی نقطہ کے لیے کون سی اصلاح استعمال ہوگی:}{\urdu{حقیقی جماعتی حدود}}{\urdu{جماعتی نشان}}{\urdu{جماعتی حدود}}{\urdu{جماعتی تعدد}}
\umcq{474}{جماعتی وقفے کا درمیانی نقطہ معلوم کرنے کا فارمولا کیا ہے:}{$Midpoint=\dfrac{Lower\ limit+Upper\ time\ limit}{2}$}{$Midpoint=\dfrac{Lower\ limit-Upper\ limit}{2}$}{$Midpoint=\dfrac{Lower\ limit\times Upper\ limit}{2}$}{None of these}
\umcq{475}{کلاس \lr{10-14} میں فریکونسی \lr{10} ہے پھر مکمل نمبر میں لکھا جاتا ہے۔:}{\lr{\tb\tb}}{\lr{\tb\tb\tb\tb\tb\tb}}{\lr{\tb\tb\tb\tb}}{None}
\umcq{476}{\lr{5-9} کی حقیقی جماعتی حدود ہے:}{5.5-10.5}{4.5-8.5}{4.5-9.5}{4-10}
\umcq{477}{پہلے \lr{20} مکمل اعداد کا عادہ ہے:}{0}{20}{19}{No Mode}
\umcq{478}{دیئے گئے مواد کا وسطانیہ کیا ہے \lr{5,8,12,13,16,18,20}:}{12}{13}{16}{18}
\umcq{479}{دیئے گئے مواد کا عادہ کیا ہے \lr{10,15,15,20,20,20,25}:}{15}{20}{25}{No mode}
\umcq{480}{کسی مواد میں بار بار آنے والی قدر کے لیے کون سی قدر استعمال ہوتی ہے:}{\urdu{اوسط}}{\urdu{وسطانیہ}}{\urdu{عادہ}}{\urdu{سعت}}
\umcq{481}{اوزانی اوسط کا فارمولا کیا ہے:}{\urdu{اوزانی اوسط}\lr{ $=\dfrac{\sum(wx)}{\sum w}$}}{\urdu{اوزانی اوسط}\lr{ $=\dfrac{\sum(wx)\times}{\mathrm{Ew}}$}}{\urdu{اوزانی اوسط}\lr{ $=\dfrac{\sum(wx)-}{\Sigma\omega}$}}{\urdu{اوزانی اوسط}\lr{ $=\dfrac{\sum(wx)+}{\mathrm{Ew}}$}}
\umcq{482}{اگر اوزان \lr{4,3,5,6} اور قدریں \lr{74,78,74,90} ہوں تو اوزانی اوسط کیا ہوگا:}{80}{85}{90}{95}
\umcq{483}{اوسط کی اصطلاح کیا ہے جہاں مختلف مشاہدات کا وزن مختلف ہوتا ہے:}{\urdu{ارتھمیٹک مین}}{\urdu{اوزانی مین}}{\urdu{میڈین}}{\urdu{موڈ}}
\umcq{484}{اگر پیکٹوں کی کل تعداد \lr{114} ہے اور علاقوں کی تعداد \lr{7} ہے تو کسی کمپنی کی طرف سے تقسیم کردہ ادویات کے پیکٹوں کی اوسط تعداد کتنی ہے:}{16}{17}{18}{19}
\umcq{485}{اگر سالانہ بجٹ \lr{5،7،8،6،7} ملین ہوں تو اگلے سال کا اوسط بجٹ کیا ہے:}{6.6}{7.0}{7.5}{8.0}
\umcq{486}{حسابی اوسط کو ظاہر کیا جاتا ہے؟:}{$\widetilde{X}$}{$\overline{X}$}{$\widehat{X}$}{SD}
\umcq{487}{اگر نمبر \lr{84,91,72,68,87,78} ہوں تو طالب کے مجموعی نمبر کتنے ہوں گے:}{}{}{}{}
\umcq{488}{جمع کے لیے استعمال ہونے والی علامت:}{$\equiv$}{$\sum$}{\%}{$\sim$}
\umcq{489}{دیئے گئے مواد کا وسطانیہ کیا ہوگا \lr{55,53,54,58,60,61,62,56,57,51,63}:}{56}{57}{58}{59}
\umcq{490}{غیر گروہی مواد کے لیے حسابی اوسط کا فارمولا استعمال ہوتا ہے:}{$\dfrac{n+1}{2}$ th term}{$\dfrac{\sum x}{n}$}{\urdu{الف اور ب دونوں}}{\urdu{ان میں سے کوئی نہیں}}
\umcq{491}{دیئے گئے مواد کا وسطانیہ کیا ہوگا \lr{110,125,122,130,124,127,120}:}{124}{125}{126}{127}
\umcq{492}{دیئے گئے مواد کا عادہ کیا ہوگا \lr{88,70,72,125,95,81,90,95,90}:}{90}{95}{\lr{90 and 95}}{No mode}
\umcq{493}{علامت \lr{$\widetilde{X}$} کس لیے استعمال ہوتی ہے:}{\urdu{عادہ}}{\urdu{وسطانیہ}}{\urdu{سعت}}{\urdu{حسابی اوسط}}
\umcq{494}{دیئے گئے مواد کا حسابی اوسط کیا ہوگا \lr{4,6,10,12,15,20,25,28,30}:}{16.67}{18.00}{20.00}{22.00}
\umcq{495}{دیئے گئے مواد کا حسابی اوسط کیا ہوگا \lr{12,18,19,0,-19,-18,-12}:}{0}{1}{2}{3}
\umcq{496}{\lr{$l+\dfrac{f_m-f_1}{(f_m-f_1)+(f_m-f_2)}\times h$} استعمال ہوتا ہے:}{\urdu{موڈ}}{\urdu{میڈین}}{\urdu{رینج}}{\urdu{حسابی اوسط}}
\umcq{497}{کون سا مواد صرف مخصوص اقدار لیتا ہے:}{\urdu{مسلسل مواد}}{\urdu{غیر مسلسل مواد}}{\urdu{گروہی مواد}}{\urdu{غیر گروہی مواد}}
\umcq{498}{مواد میں کسی قدر کے وقوع پذیر ہونے کی تعداد کو کہا جاتا ہے:}{\urdu{تعدد}}{\urdu{متوقع تعدد}}{\urdu{جماعتی حد}}{\urdu{حقیقی جماعتی حد}}
\umcq{499}{درمیانی نقطہ جانا جاتا ہے:}{\urdu{حسابی اوسط}}{\urdu{وسطانیہ}}{\urdu{جماعتی حد}}{\urdu{جماعتی نشان}}
\umcq{500}{تعددی کثیر الاضلاع کو استعمال کرتے ہوئے بنایا جاتا ہے:}{\urdu{کالمی نقشہ}}{\urdu{بار گراف}}{\urdu{حقیقی جماعتی حدود}}{\urdu{سعت}}
\umcq{501}{سب سے بڑی مد اور سب سے چھوٹی مد کے درمیان فرق کو کہا جاتا ہے:}{\urdu{جماعتی حدود}}{\urdu{درمیانی نقطہ}}{\urdu{نسبتی تعدد}}{\urdu{سعت}}
\umcq{502}{مرکزی رجحان کی پیمائش معلوم کرنے کے لیے استعمال کیا جاتا ہے:}{\urdu{حقیقی جماعتی حدود}}{\urdu{مجموعی حدود}}{\urdu{درمیانی قدر}}{\urdu{تعدد}}

\umcq{503}{اگر \lr{5,7,8,9} اور \lr{x} کا حسابی اوسط \lr{7,5} ہو تو \lr{x} کی قدر کیا ہوگی:}{10}{8}{8.5}{5.8}
\umcq{504}{دیے گئے مواد کا عادہ معلوم کریں \lr{2,5,8,9,0,1,3,7} اور \lr{10}:}{5}{7}{0}{No mode}
\umcq{505}{کسی سلسلہ یا مواد میں وہ قیمت جو سب سے زیادہ مرتبہ آئے کہلاتی ہے:}{\urdu{حسابی اوسط}}{\urdu{عادہ}}{\urdu{وسطانیہ}}{\urdu{اوزانی اوسط}}
\umcq{506}{دیے گئے مواد کا وسطانیہ معلوم کریں \lr{110,125,122,130,124} اور \lr{120}:}{}{}{}{}
\umcq{507}{احتمال کا فارمولا کیا ہے:}{$P(A)=\dfrac{n(A)}{n(S)}$}{$P(A)=\dfrac{n(S)}{n(A)}$}{P(A) = n(A) x n(S)}{P(A) = n(S) $-$ n(A)}
\umcq{508}{درج ذیل میں سے کونسی تجربے کی مثال نہیں ہے:}{\urdu{سکہ ٹاس کرنا}}{\urdu{ڈائس کا پھینکنا}}{\urdu{ایک نمبر کا اندازہ لگانا}}{\urdu{تاش میں سے ایک پتا نکالنا}}
\umcq{509}{جب ایک سکے کا اچھالا جاتا ہے تو ہیڈ آنے کا احتمال کیا ہوگا:}{$\dfrac{1}{2}$}{$\dfrac{1}{3}$}{$\dfrac{1}{4}$}{$\dfrac{1}{6}$}
\umcq{510}{اگر ایک ڈائس پھینکا جائے تو \lr{3} پر تقسیم ہونے والے نمبر آنے کا احتمال کیا ہوگا:}{$\dfrac{1}{6}$}{$\dfrac{1}{4}$}{$\dfrac{1}{3}$}{$\dfrac{1}{2}$}
\umcq{511}{جب ایک سنگل سکے کا اچھالا جاتا ہے تو سیمپل سپیس کیا ہوگا:}{\{A, B\}}{\{1, 2, 3\}}{\{1, 2, 3, 4\}}{\{H, T\}}
\umcq{512}{اگر ایک ڈائس کو دو دفعہ پھینکا جائے تو ممکنہ نتائجے کل نمبر کیا ہوں گے:}{6}{36}{12}{24}
\umcq{513}{اگر ایک واقعہ \lr{A} کے \lr{3} پسندیدہ نتائج ہوں۔ اور سیمپل سپیس \lr{6} ممکن نتائج رکھتی ہو تو \lr{P(A)} کیا ہوگا:}{$\dfrac{1}{6}$}{$\dfrac{1}{4}$}{$\dfrac{1}{3}$}{$\dfrac{1}{2}$}
\umcq{514}{ایونٹ \lr{A} کی تکمیل کیا ہوگی:}{P(S)}{P(A)}{P(A$'$)}{P(A and B)}
\umcq{515}{ایک پھینکے گئے ڈائس پر \lr{6} نہ آنے کا احتمال کیا ہوگا:}{$\dfrac{1}{6}$}{$\dfrac{5}{6}$}{$\dfrac{1}{3}$}{$\dfrac{2}{3}$}
\umcq{516}{اگر ایک ڈائس کا پھینکا جائے تو ایک جفت عدد آنے کا احتمال کیا ہوگا:}{$\dfrac{1}{2}$}{$\dfrac{1}{3}$}{$\dfrac{1}{4}$}{$\dfrac{1}{6}$}
\umcq{517}{جب ایک ڈائس کو دو دفعہ پھینکا جاتا ہے تو نمبر \lr{8} کا مجموعہ آنے کا احتمال کیا ہوگا:}{$\dfrac{1}{9}$}{$\dfrac{5}{36}$}{$\dfrac{1}{36}$}{$\dfrac{1}{12}$}
\umcq{518}{جب دو ڈائسوں کو پھینکا جائے تو ڈبل چھ آنے کا احتمال کیا ہوگا:}{$\dfrac{1}{36}$}{$\dfrac{1}{18}$}{$\dfrac{1}{12}$}{$\dfrac{1}{6}$}
\umcq{519}{کسی بھی ایونٹ کا احتمال ہمیشہ ہوتا ہے:}{\urdu{صفر سے کم}}{\urdu{\lr{1} سے زیادہ}}{\urdu{\lr{1} کے برابر}}{\urdu{\lr{0} اور \lr{1} کے درمیان}}
\umcq{520}{:}{$\dfrac{5}{13}$}{$\dfrac{8}{13}$}{$\dfrac{1}{3}$}{$\dfrac{5}{8}$}
\umcq{521}{\lr{52} کے ڈیک سے کارڈ کھینچتے وقت نہ تو اسپیڈ اور نہ ہی دل ملنے کا کیا امکان ہے:}{$\dfrac{26}{52}$}{$\dfrac{26}{39}$}{$\dfrac{1}{2}$}{$\dfrac{1}{4}$}
\umcq{522}{اگر کسی ایونٹ کا احتمال \lr{0.3} ہو تو اس کی تکمیل کا امکان ہوگا:}{0.4}{0.5}{0.3}{0.7}
\umcq{523}{تاش کے \lr{52} پتوں میں سے سرخ رنگ کا دل نکالنے کا احتمال کیا ہوگا:}{$\dfrac{1}{13}$}{$\dfrac{1}{4}$}{$\dfrac{1}{2}$}{$\dfrac{1}{26}$}
\umcq{524}{ایک بستہ جس میں سرخ اور \lr{5} نیلی گیندیں ہوں اس میں سے سرخ گیند نہ آنے کا احتمال کیا ہوگا:}{$\dfrac{5}{13}$}{$\dfrac{8}{13}$}{$\dfrac{5}{8}$}{$\dfrac{8}{5}$}
\umcq{525}{اگر \lr{3} ڈائسوں کو گھمایا جائے تو کم از کم ایک \lr{6} آنے کا احتمال کیا ہوگا:}{$\dfrac{1}{36}$}{$\dfrac{3}{36}$}{$\dfrac{1}{6}$}{$\dfrac{5}{36}$}
\umcq{526}{جب دو ڈائسوں کو گھمایا جائے تو جفت مجموعہ آنے کا احتمال کیا ہوگا:}{$\dfrac{5}{36}$}{$\dfrac{18}{36}$}{$\dfrac{9}{36}$}{$\dfrac{12}{36}$}
\umcq{527}{تاش کے \lr{52} پتوں میں سے ایک سرخ کارڈ کا احتمال کیا ہوگا:}{$\dfrac{1}{2}$}{$\dfrac{13}{52}$}{$\dfrac{26}{52}$}{$\dfrac{1}{26}$}
\umcq{528}{ایک بستہ جس میں \lr{5} سفید اور \lr{7} کالے گیند موجود ہوں تو اس میں ایک سفید گیند حاصل کرنے کا احتمال کیا ہوگا:}{$\dfrac{1}{3}$}{$\dfrac{7}{12}$}{$\dfrac{1}{2}$}{$\dfrac{5}{12}$}

\umcq{529}{نسبتی تعدد ہے:}{\urdu{پسندیدہ نتائج والے نمبرز}}{\urdu{ممکنہ نتائج کے نمبرز}}{\urdu{ایک ایونٹ اکثر کیسے رونما ہوتا ہے}}{\urdu{ٹرائلز کے نمبرز}}
\umcq{530}{\lr{80} طلبہ میں سے ایک سرخ گیند کا انتخاب کرنے کا نسبتی تعدد کیا ہوگا:}{0.29}{0.19}{0.31}{0.12}
\umcq{531}{اگر چھ ڈائس کو \lr{50} بار گھمایا جاتا ہے، تو ظاہر ہونے والی تعداد کی متوقع تعدد اس کے ذریعہ دی جاتی ہے:}{NP(B)}{NP(A)}{NP(X)}{None of these}
\umcq{532}{\lr{80} طلبہ کے سروے میں سے سبز کا احتمال کیا ہوگا؟؟:}{$\dfrac{15}{80}$}{0.19}{0.31}{$\dfrac{10}{80}$}
\umcq{533}{اگر احتمال \lr{$\dfrac{1}{5}$} ہو تو، \lr{100} ٹرائلز میں متوقع تعدد کیا ہوگا:}{20}{50}{10}{25}
\umcq{534}{سیمپل اسپیس کا ہر رکن کہلاتا ہے:}{\urdu{واقعہ}}{\urdu{تجربہ}}{\urdu{سیمپل نقطہ}}{\urdu{نتائج}}
\umcq{535}{ایک نتیجہ جو اس بات کو ظاہر کرتا ہے کہ ہم کتنی بار واقعہ کے وقوع پذیر ہونے کی توقع کرتے ہیں کہلاتا ہے:}{\urdu{نتائج}}{\urdu{موافق نتیجہ}}{\urdu{سیمپل اسپیس}}{\urdu{سیمپل نقطہ}}
\umcq{536}{ہمیں کس سے معلوم ہوتا ہے کہ ایک مخصوص واقعہ گل تعدد کے مقابلے میں کتنی بار آتا ہے:}{\urdu{متوقع تعدد}}{\urdu{نسبتی تعدد کا مجموعہ}}{\urdu{نسبتی تعدد}}{\urdu{تعدد}}
\umcq{537}{کسی واقعہ کے وقوع پذیر ہونے کا متوقع احتمال اس طرح بھی جانا جاتا ہے:}{\urdu{نسبتی تعدد}}{\urdu{متوقع تعدد}}{\urdu{حقیقی جماعتی حدود}}{\urdu{متوقع تعدد کا مجموعہ}}
\umcq{538}{تمام متوقع تعدد کا مجموعہ کس مقررہ عدد کے برابر ہے:}{\urdu{کل تجربات}}{\urdu{متوقع تعدد}}{\urdu{نتائج}}{\urdu{واقعات}}
\umcq{539}{جب چار ڈائس پھینکنے جاتے ہیں تو سیمپل اسپیس کی کل تعداد معلوم کریں:}{6$^{2}$}{6$^{3}$}{6$^{4}$}{6$^{6}$}
\umcq{540}{کسی خاص واقعہ کے وقوع پذیر ہونے کا امکان کہلاتا ہے:}{\urdu{سیمپل سپیس}}{\urdu{متوقع احتمال}}{\urdu{احتمال}}{\urdu{متوقع تعدد}}
\umcq{541}{اگر ڈائس کا جوڑا پھینکا جائے تو دوبار \lr{2} آنے کا احتمال کیا ہوگا:}{$\dfrac{1}{6}$}{$\dfrac{1}{3}$}{$\dfrac{5}{6}$}{$\dfrac{1}{36}$}
\umcq{542}{تاش کے \lr{52} پتوں میں سے ایک پتہ منتخب کیا، غلام اور بادشاہ نا آنے کا کیا احتمال ہوگا:}{$\dfrac{2}{13}$}{$\dfrac{11}{13}$}{$\dfrac{2}{52}$}{$\dfrac{11}{52}$}
\umcq{543}{ایسا واقعہ جو رونما ہوگا اور اس کے ہونے کا امکان زیادہ ہو کہلاتا ہے:}{\urdu{مساوی ممکنہ واقعہ}}{\urdu{ممکنہ واقعہ}}{\urdu{یقینی واقعہ}}{\urdu{غیر ممکنہ واقعہ}}

\end{document}
